\section{\PNR: revenue-funded operator compensation}
\label{sec:pnr}\label{sec:07-pnr}

Everything a tenant pays to run an application on \Flux today lands at a single transparent address,
and every satoshi a node operator earns is newly issued \shipped
\srcmeasured{api.runonflux.io/\allowbreak{}apps/\allowbreak{}deploymentinformation}{2026-09-03},
\srccode{fluxd/\allowbreak{}src/\allowbreak{}fluxnode/\allowbreak{}fluxnode.cpp}{892--936}. Progressive Node Rewards
(\PNR) is the planned change to the second half of that sentence, and it is designed so that the
first half would look unchanged to whoever is paying: the tenant would still make \emph{one}
payment, of the full price, to \emph{one} address, and consensus --- not the payer, not the wallet,
and not \FluxOS --- would divide what arrives there. A fixed fraction would accrue to a
protocol-held fee pool, and that pool would be paid out to node operators through the block
subsidy's existing payee rotation, one small increment per block \planned
\srcintent{evidence/design/08-answers-round-2.md}. The purpose of this section is to state that
mechanism precisely enough to implement, to
prove the fairness property the design rests on, to prove that the pooling defeats a timing attack
that would otherwise invert the mechanism, and --- with equal care --- to state the two things the
mechanism does \emph{not} do: it does not make operator income proportional to collateral, and it
creates no incentive whatsoever to host anything.

\paragraph{Terminology hazard, stated once and before anything else.}
\PoN (\emph{Proof of Node}) and \PNR (\emph{Progressive Node Rewards}) are different mechanisms, in
different layers, at different maturities, and the names are trivially confused. \PoN is the
deployed consensus rule that replaced hashing: it selects one payee per tier per block from a
collateral-gated queue and pays that payee out of the block subsidy \shipped
\srccode{fluxd/\allowbreak{}src/\allowbreak{}fluxnode/\allowbreak{}fluxnode.cpp}{713--778}. \PNR is the planned revenue-sharing rule
specified in this section: it takes a fraction of application payments and routes it to whichever
nodes \PoN has already selected \planned \srcintent{evidence/design/02-pnr.md}. \PNR consumes \PoN's
payee set and adds no selection mechanism of its own. Wherever this section says ``the rotation'' it
means \PoN's rotation; wherever it says ``the pool'' or ``the drip'' it means \PNR.

\paragraph{Symbols local to this section.}
Beyond \texttt{notation.tex} this section uses: $F_\hgt$ for the pool inflow of block $\hgt$ (in
$\sat$); $G_\hgt$ for the Foundation share of the same block's receipts; $D_\hgt$ for the drip paid
out of the coinbase of block $\hgt$ and $D_{\hgt,t}$ for its tier-$t$ part; $\epsilon_\hgt$ for the
integer partition residue and $U_\hgt$ for tier parts that find no payee; $A_F$ and $A_\Phi$ for the
Foundation and pool parts into which consensus divides a single received payment; $\mathsf{sink}$
for the consensus application-payment address, $\mathcal{O}_\hgt$ for the set of outputs paying it in
block $\hgt$ and $\mathrm{Rev}_\hgt$ for its balance at height $\hgt$; $m_\app$ for the canonical
serialisation of application $\app$'s specification and $H$ for the hash carried in the payment memo;
$q \defeq 1 - 1/\drip$; $r(\hgt) \defeq 100\,\nodeshare(\hgt)$ for the integer form of the operator
share; $V$ for annual application revenue in $\flux$ per year; $g$ for annual revenue growth; $m$
for an exogenous price multiple used only in the clearly-labelled secondary analysis of
\S\ref{sec:pnr:real}; and $w_{\adv}$ for an adversary's share of the drip. All consensus arithmetic
is over integers in $\sat$, with $10^{8}\sat = 1\flux$
\srccode{fluxd/\allowbreak{}src/\allowbreak{}amount.h}{17}.

% ---------------------------------------------------------------------------
\subsection{Assumptions}
\label{sec:pnr:assumptions}

\begin{assumption}[Model facts this section depends on]
\label{as:pnr}
All properties below are stated against the system and adversary model of \S2. Five of its facts are
load-bearing here and are restated so that no proof silently assumes more than the model gives.
\begin{enumerate}
  \item[\textbf{A1}] \emph{Deterministic payee selection.} For each tier $t \in \tiers$ the payee of
    block $\hgt$ is the front of a queue sorted by last-paid height ascending, with never-paid nodes
    keyed by confirmation height; the front entry is returned and stale entries purged \shipped
    \srccode{fluxd/\allowbreak{}src/\allowbreak{}fluxnode/\allowbreak{}fluxnode.h}{313--328},
    \srccode{fluxd/\allowbreak{}src/\allowbreak{}fluxnode/\allowbreak{}fluxnode.cpp}{713--778}. The selection is a public, deterministic
    function of chain state --- it is predictable to every observer, and \S\ref{sec:pnr:timing}
    treats that predictability as the adversary's capability, not as an oversight.
  \item[\textbf{A2}] \emph{Fixed chainwork.} Every \PoN block contributes a fixed $2^{40}$ to
    chainwork, so fork choice is a block-count race rather than a cumulative-work comparison
    \shipped \srccode{fluxd/\allowbreak{}src/\allowbreak{}pow.cpp}{284--290} (C14). No proof below may appeal to reorg cost
    scaling with hashing.
  \item[\textbf{A3}] \emph{Emergency block production exists.} A hardcoded 2-of-4 key set can
    produce a block outside the normal \PoN path, with no time or stall-detection gating, and the
    remainder output of its coinbase defaults to the dev-fund address while its tier outputs remain
    bound by the deterministic payout check \shipped
    \srccode{fluxd/\allowbreak{}src/\allowbreak{}emergencyblock.cpp}{183--191},
    \srccode{fluxd/\allowbreak{}src/\allowbreak{}rpc/\allowbreak{}emergencyblock.cpp}{64--77},
    \srccode{fluxd/\allowbreak{}src/\allowbreak{}miner.cpp}{487--489} (C9). This is an adversary capability the
    mechanism must be designed against, not a hypothetical.
  \item[\textbf{A4}] \emph{The eligibility gate is not operator-bound and not machine-bound.} The
    benchmark signature that admits a node to the payee set is produced by a process on the
    operator's own host under network-wide shared keys (C13), and the \ArcaneOS attestation
    signature is derived from values compiled identically into every binary (C22). Anything \PNR
    pays for is gated by exactly this, and by nothing stronger.
  \item[\textbf{A5}] \emph{Consensus can recognise only what it can parse --- and a payment to a
    fixed address is parseable.} Consensus can identify value received at an address it carries as a
    constant, and can therefore divide that value without the payer's cooperation and without any
    off-chain party's. What consensus cannot decide is whether a given application \emph{ought} to
    have been paid for: the admission rule that binds a payment to a deployment is a \FluxOS
    validation rule, and the \flux-denominated price table is Foundation policy. The precise scope
    of ``consensus-enforced'' is thus: \emph{the distribution of application revenue is a consensus
    rule; the obligation to pay, and the price, are policy} \srcdoc{plan/conflicts.md C38}.
\end{enumerate}
\end{assumption}

\begin{assumption}[Fleet and revenue baseline]
\label{as:pnr:fleet}
Fleet figures are $\ncount{\tC} = 3{,}247$, $\ncount{\tN} = 1{,}615$, $\ncount{\tS} = 1{,}627$,
total $6{,}489$ \srcmeasured{api.runonflux.io/\allowbreak{}daemon/\allowbreak{}listzelnodes}{2026-09-03}. The settlement model
and its generated data files were computed at $\ncount{\tC} = 3{,}246$
\srcdoc{scripts/07-pnr-settlement.py, run 2026-09-03}; the one-node difference changes nothing at
the precision reported. The activation-scale revenue figure $V_0 = 1{,}431{,}600 \flux$ per year is
the team's assumption of $1{,}193$ applications at $100\flux$ per month, \emph{not} a measured
receipt total \srcintent{evidence/design/02-pnr.md}; the census over the same window returned
$1{,}195$ application records and $7{,}486$ requested instances
\srcmeasured{api.runonflux.io/\allowbreak{}apps/\allowbreak{}globalappsspecifications}{2026-09-03}. Every absolute
\flux-denominated figure in \S\ref{sec:pnr:crossover} is linear in $V_0$; the growth-rate
conclusions are not.
\end{assumption}

% ---------------------------------------------------------------------------
\subsection{The payment, and where the split is applied}
\label{sec:pnr:tx}

Today an application is charged for by a transparent payment to a single policy-configured address,
\texttt{t3Nryf\allowbreak{}AQLGe\allowbreak{}Fs9j\allowbreak{}Eoeqsxm\allowbreak{}BN2QLRa\allowbreak{}RKFLUX}; \FluxOS checks that the payment exists and covers the
specification, and \fluxd sees an ordinary transfer \shipped
\srcmeasured{api.runonflux.io/\allowbreak{}apps/\allowbreak{}deploymentinformation}{2026-09-03},
\srccode{flux/\allowbreak{}ZelBack/\allowbreak{}src/\allowbreak{}services/\allowbreak{}appMessaging/\allowbreak{}messageVerifier.js}{734--763} (C9). The carrier that binds a
payment to what it buys is also already in the fleet: \CumulusVPN pays a fixed address with an
\texttt{OP\_RETURN} memo, and every gateway independently derives entitlement by scanning the chain
for correctly-memoed payments, with no account and no server-issued credential \shipped
\srccode{cumulusvpn/\allowbreak{}gateway/\allowbreak{}internal/\allowbreak{}entitle/\allowbreak{}entitle.go}{99--108,218--287},
\srccode{cumulusvpn/\allowbreak{}clients/\allowbreak{}core-ts/\allowbreak{}src/\allowbreak{}paymentCode.ts}{17--23} (C19).

\PNR would keep that shape at the payer and move the address into consensus. One output, one memo,
no second recipient, no fee component and no new transaction type: the settled decision is that the
tenant would pay the full price once, and the division would be performed by rules in which no
wallet, no agent and no \FluxOS release participates \planned
\srcintent{evidence/design/08-answers-round-2.md} \srcdoc{plan/conflicts.md C36}.

\begin{definition}[\PNR application payment, \planned]
\label{def:pnr:tx}
Fix the operator-share schedule
\[
  \mathrm{red}(\hgt) \defeq \min\!\left\{\redmax,\ \floorfrac{\hgt - \hpon}{\redint}\right\},
  \qquad
  \nodeshare(\hgt) \defeq
  \begin{cases}
    0, & \hgt < \hpa,\\[2pt]
    \min\!\big\{\nodesharecap,\ 0.20 + \rampstep\,(\mathrm{red}(\hgt) - 1)\big\}, & \hgt \ge \hpa,
  \end{cases}
\]
with $\hpon = 2{,}020{,}000$ \srccode{fluxd/\allowbreak{}src/\allowbreak{}chainparams.cpp}{153},
$\redint = 1{,}051{,}200$ blocks \srccode{fluxd/\allowbreak{}src/\allowbreak{}chainparams.cpp}{157},
$\redmax = 20$ \srccode{fluxd/\allowbreak{}src/\allowbreak{}chainparams.cpp}{158},
$\hpa = 3{,}466{,}630$, $\rampstep = 0.05$ and $\nodesharecap = 0.50$
\srcintent{evidence/design/02-pnr.md}. Write $r(\hgt) \defeq 100\,\nodeshare(\hgt) \in \mathbb{Z}$.

Let $\mathsf{sink}$ denote the \emph{application-payment address}: one transparent address carried
as a consensus constant from height $\hpa$ onward, whose identity is O14. Let $m_\app$ be the
canonical serialisation of the specification of application $\app$ and $H$ the hash already used for
application identity, so $H(m_\app) \in \{0,1\}^{256}$.

A \emph{\PNR application payment} for $\app$ at height $\hgt \ge \hpa$ is an ordinary transparent
transaction of the shape PS1--PS2, over an address governed by the consensus rules PS3--PS4. The
division of the four is itself the content of the design and is not cosmetic: PS1--PS2 are what a
payer constructs and what \FluxOS checks, PS3--PS4 are what \fluxd applies to \emph{everything} that
arrives at the address, memo or no memo.
\begin{description}
  \item[PS1 (one recipient, full price).] The transaction carries exactly one value-bearing output
    paying $\mathsf{sink}$, of value $v \in \mathbb{Z}_{\ge 0}$ in $\sat$, with
    $v \ge \price(\app,\hgt)$ (Definition~\ref{def:pnr:price}). The payer constructs one recipient
    and its own change, and performs no other arithmetic.
  \item[PS2 (memo).] The transaction carries one zero-value \texttt{OP\_RETURN} output whose data
    field is the \num{32}-byte value $H(m_\app)$. Consensus carries the memo and does not interpret
    it; \FluxOS reads it to attribute the payment to an application, and rejects a registration
    whose payment is missing, mis-memoed or short of $\price(\app,\hgt)$. A sink payment carrying no
    memo is still split by PS4 --- it simply buys nothing.
  \item[PS3 (the sink is unspendable).] No transaction may spend an output paying $\mathsf{sink}$;
    a block containing such a spend is invalid. Value paid to $\mathsf{sink}$ leaves circulation at
    the moment of payment.
  \item[PS4 (consensus split).] On connecting the block containing the payment at height $\hgt$,
    consensus divides the received value: $A_\Phi \defeq \floorfrac{r(\hgt)\,v}{100}$ is credited to
    the fee pool $\pool$, and $A_F \defeq v - A_\Phi$ is credited to the Foundation and paid through
    the \emph{existing} dev-fund coinbase output. The floor is taken once, over integers, per
    received output, and favours the Foundation by strictly less than $1\sat$ per payment.
\end{description}
Write $F_\hgt$ and $G_\hgt$ for the sums of $A_\Phi$ and of $A_F$ over the outputs paying
$\mathsf{sink}$ in block $\hgt$; gross receipts in that block are $F_\hgt + G_\hgt$, and
$\mathrm{Rev}_\hgt \defeq \sum_{j = \hpa}^{\hgt}\big(F_j + G_j\big)$ is the balance held at
$\mathsf{sink}$. The domain of PS3 and PS4 is $\hgt \ge \hpa$; below the activation height a payment
to the deployed sink is an ordinary spendable output, so if O14 retains the address in use today,
$\mathrm{Rev}$ is defined from $\hpa$ and not from that address's first receipt.
\end{definition}

The three decisions inside Definition~\ref{def:pnr:tx} that were open in the first draft of this
specification are now settled and are recorded here rather than left as parameters: the memo binds
$H(m_\app)$ and is carried in an \texttt{OP\_RETURN}, which closes O3; the payment is not a new
transaction type and needs no new serialisation, which closes half of O18; and the split is applied
by consensus rather than by the payer or by \FluxOS, which closes the other half
\srcintent{evidence/design/08-answers-round-2.md} \srcdoc{plan/conflicts.md C36, C38}.

\begin{remark}[O7 --- resolved by the retirement of the shielded pools]
\label{rem:o7-moot}
PS4 is computable only if the value received at $\mathsf{sink}$ is visible to consensus, and the
question was which of three regimes applied to shielded payers. It is answered by a decision taken
elsewhere: both shielded pools are to be retired (\S\ref{sec:shielded-retirement}), so
\textbf{all settlement is transparent} and PS4 is computable unconditionally. The domain collapses to
its first element, not by argument within this section but because the alternative ceases to exist.

Worth noting for the record: an intermediate answer had permitted shielded payers with a
revealed-amount output \srcintent{08-answers-round-2.md, round 9}, which would have made PS4
computable while hiding the payer. That is a coherent design and it is now moot. \PNR{} is
consequently simpler than any draft of this section assumed --- every receipt at $\mathsf{sink}$ is
publicly attributable, and the fee pool's inflow is auditable by anyone.
\end{remark}

The coinbase already carries four value-bearing outputs that consensus checks: a dev-fund minimum
paid to a hardcoded address, and one output per tier paid to that tier's \PoN payee \shipped
\srccode{fluxd/\allowbreak{}src/\allowbreak{}main.cpp}{2877--2884}, \srccode{fluxd/\allowbreak{}src/\allowbreak{}main.cpp}{1231--1249},
\srccode{fluxd/\allowbreak{}src/\allowbreak{}fluxnode/\allowbreak{}fluxnode.cpp}{837--875} (C9). Those four outputs would be the whole of
the machinery \PNR needs on the paying side.

\begin{property}[The split is a consensus rule, and the single output is why, \planned]
\label{prop:pnr:enforce}
Under Definition~\ref{def:pnr:tx} the \SI{80}{\percent}/\SI{20}{\percent} division would be applied
by \fluxd alone. Three consequences follow, and they are stated together because they follow from
one decision.
\begin{enumerate}
  \item[\textbf{E1}] \emph{The payer could not get it wrong.} There would be one recipient and one
    amount, so there would be no split for a wallet to compute, mis-round or omit, and no second
    recipient for a wallet that exposes only one to refuse. Any wallet that can pay an address and
    attach an \texttt{OP\_RETURN}, and any agent holding a bare wallet API, could construct a valid
    payment --- which is the property \S18's agent-native argument rests on, and which the deployed
    \CumulusVPN payment path cited above is the existence proof for.
  \item[\textbf{E2}] \emph{\FluxOS could not change it.} \FluxOS would check only that a payment
    exists, carries $H(m_\app)$ and covers $\price(\app,\hgt)$. The fraction $\nodeshare(\hgt)$, the
    tier partition $\tierratio{t}$ and the destination of either share would be alterable only by a
    consensus change, never by a \FluxOS release through the channel of \S12 and \S14.
  \item[\textbf{E3}] \emph{The coinbase would gain no output.} The Foundation share would be added
    to the dev-fund output and the drip to the three tier outputs, and nothing else would be added.
    The team's constraint that the coinbase must not grow \srcintent{evidence/design/02-pnr.md}
    would be met strictly, and met better than by any shape that adds an output.
\end{enumerate}
\end{property}

\begin{proof}[Justification]
E1 is immediate from PS1: the transaction has one $\mathsf{sink}$ output and the payer's change, and
PS4 is evaluated by the validator, not by the payer. E3 is immediate from PS4 and from
Definition~\ref{def:pnr:pool}, which route $G_\hgt$ and $D_{\hgt,t}$ into outputs that already
exist. E2 is the substantive one and is a statement about \emph{where a rule can be edited}. A
\FluxOS release can change any predicate \FluxOS evaluates; it cannot change a predicate evaluated
in \texttt{ConnectBlock}, because a node running the changed rule would be rejected by every other
node. Under Definition~\ref{def:pnr:tx} the division is evaluated in \texttt{ConnectBlock} and the
admission rule is evaluated in \FluxOS, so the two are separated by the strongest boundary the
system has. Note precisely what E2 does not say: it does not say the \emph{amount} is beyond policy.
The price table is Foundation-set (Definition~\ref{def:pnr:price}, O9) and which channels settle
through $\mathsf{sink}$ at all is a \FluxOS rule (O8). Consensus would fix the distribution; policy
would fix the base.
\end{proof}

\begin{property}[Supply neutrality, and a public revenue counter, \planned]
\label{prop:pnr:supply}
Let $\supply(\hgt)$ be the issuance function of \S5, and let $\supply'(\hgt)$ denote circulating
supply under \PNR --- issuance, less value locked at $\mathsf{sink}$ by PS3, plus the coinbase
re-issuance of PS4 and of the drip. Then for all $\hgt \ge \hpa$,
\[
  \supply'(\hgt) \;=\; \supply(\hgt) \;-\; \pool_\hgt .
\]
\PNR would therefore create no coin and destroy none; the only deviation from \S5's schedule would
be the pool balance in flight, which is bounded by $\drip F_{\max}$
(Property~\ref{prop:pnr:invariants}, I4) and equals $\SI{30}{\day}$ of the operator share at steady
state. Separately, $\mathrm{Rev}_\hgt$ --- the balance of a single address --- would be the
cumulative gross application revenue settled through this path since $\hpa$.
\end{property}

\begin{proof}
By PS3 the value at $\mathsf{sink}$, namely
$\mathrm{Rev}_\hgt = \mathrm{In}_\hgt + \sum_{j \le \hgt} G_j$ with
$\mathrm{In}_\hgt = \sum_{j \le \hgt}F_j$, is unspendable, so circulating supply is reduced by
exactly that amount. It re-enters through the coinbase in two streams: the Foundation share
$\sum_{j \le \hgt} G_j$ through the dev-fund output, and the operator share as the drip actually
placed in tier outputs, $\mathrm{Paid}_\hgt$ (Property~\ref{prop:pnr:invariants}). Hence
$\supply' - \supply = \sum_{j \le \hgt}G_j + \mathrm{Paid}_\hgt - \mathrm{Rev}_\hgt
 = \mathrm{Paid}_\hgt - \mathrm{In}_\hgt = -\pool_\hgt$, the last step by I1. For the counter: by
PS1 every settled payment pays $\mathsf{sink}$ and by PS3 nothing ever leaves it, so its balance is
the running sum of received values.
\end{proof}

\begin{remark}[What the counter is worth, and what it does not cover]
\label{rem:pnr:counter}
\S5 argues that supply should be verifiable by arithmetic rather than by trust. A
consensus-unspendable sink turns the application-revenue term of that arithmetic into a single
number anyone can query from a block explorer, with no party privileged to publish it and none able
to withhold it. The consequence for \S\ref{sec:evaluation} is concrete: that section presently reasons from an
\emph{assumed} $100\flux$ per application per month (Assumption~\ref{as:pnr:fleet})
\srcintent{evidence/design/02-pnr.md}, and after activation the same quantity would be measured
rather than assumed, which is also what would make O17 closable by observation. One qualification
travels with the counter and must not be dropped: it would measure what settles through this path,
so every channel excluded by O8 --- credits (\S16), \CumulusVPN, enterprise contracts --- is revenue
the counter would not see and the consensus split would not touch.
\end{remark}

\begin{definition}[Payment function, interface, \planned]
\label{def:pnr:price}
For an application $\app$ with resource vector
$\resvec = (\mathrm{cpu}, \mathrm{ram}, \mathrm{hdd}) \in \mathbb{R}_{\ge 0}^{3}$, instance count
$\ninst$, and duration $\dur \in [1,\ 1{,}056{,}000]$ blocks, the price
is a function
$\price\colon \mathcal{A} \times \mathbb{Z}_{\ge 0} \to \mathbb{Z}_{\ge 0}$, where $\mathcal{A}$
denotes the set of well-formed \appspec{}s, of the form
\[
  \price(\app, \hgt) \;=\; \max\!\Big\{\,p_{\min},\ \
    \ninst \cdot \frac{\dur}{L^{\ast}(\hgt)} \cdot
    \big(\pricevec_{\hgt}^{\!\top}\resvec + \mathrm{surcharges}(\app)\big)\Big\},
\]
where $\pricevec_\hgt$ is the height-scheduled price vector and $p_{\min} = 0.01\flux$. The
duration denominator carries a fork branch: $L^{\ast}(\hgt) = \blockslasting = 22{,}000$ blocks before
\PoN activation and $4\blockslasting = 88{,}000$ blocks from height $2{,}020{,}000$, because \PoN
quadrupled the block rate and the multiplier preserves wall-clock duration \shipped
\srccode{flux/\allowbreak{}ZelBack/\allowbreak{}src/\allowbreak{}services/\allowbreak{}appRequirements/\allowbreak{}appValidator.js}{864-869}, \srccode{fluxd/\allowbreak{}src/\allowbreak{}chainparams.cpp}{101,105}.
The instance count is bounded above at $100$; its lower bound is $3$ in general but $1$ for
\specv{8} applications from height $2{,}176{,}519$, which is the value consensus-side validation
enforces even though the published \texttt{deploymentinformation} endpoint continues to report $3$
\shipped \srccode{flux/\allowbreak{}ZelBack/\allowbreak{}src/\allowbreak{}services/\allowbreak{}appRequirements/\allowbreak{}appValidator.js}{818-828}, \srcdoc{plan/conflicts.md C33}. The duration floor is likewise $1$ block from height $1{,}964{,}447$ by the same validator, where the endpoint still reports $5{,}000$ \shipped \srccode{flux/\allowbreak{}ZelBack/\allowbreak{}src/\allowbreak{}services/\allowbreak{}appRequirements/\allowbreak{}appValidator.js}{852-863}, \srccode{flux/\allowbreak{}ZelBack/\allowbreak{}src/\allowbreak{}services/\allowbreak{}appQuery/\allowbreak{}deploymentInfoService.js}{41-46}. The deployed schedule has six rows, i.e.\ five revisions since genesis; the
current row, effective from height $1{,}597{,}156$, is
$\pricevec = (0.03,\ 0.01,\ 0.004)\flux$ per unit with surcharges
$(\text{port},\ \text{scope},\ \text{staticip}) = (0.4,\ 0.8,\ 0.4)\flux$ \shipped
\srcmeasured{api.runonflux.io/\allowbreak{}apps/\allowbreak{}deploymentinformation}{2026-09-03}. Settlement consumes
$\price$ and never recomputes it.
\end{definition}

\notestablished{the exact \FluxOS aggregation --- the units of the price table, the rounding, and
any enterprise or geolocation multipliers --- is not established by the evidence available to this
section. The price \emph{table} above is measured; the formula that consumes it is stated as an
interface and must be cited from the \FluxOS evidence before publication.}

Two consequences of Definition~\ref{def:pnr:price} deserve to be visible rather than discovered.
First, the pool would receive $\nodeshare(\hgt)\,v$ regardless of how $\price$ was computed:
consensus would divide whatever arrives at $\mathsf{sink}$, not check that what arrives is correct.
Second, any change to the price table would change pool inflow one-for-one with no consensus
involvement. The table is discretionary today and has been revised five times, the CPU rate falling
$100\times$ from $3$ to $0.03\flux$ per unit \shipped
\srcmeasured{api.runonflux.io/\allowbreak{}apps/\allowbreak{}deploymentinformation}{2026-09-03}. \emph{Consensus would fix the
distribution; policy would fix the base.} That sentence is the honest form of the decentralization
claim this mechanism supports --- it is now stronger on the first half than the first draft of this
specification could claim (Property~\ref{prop:pnr:enforce}, E2) and unchanged on the second, and
\S\ref{sec:pnr:openparams} records the two discretionary levers it leaves standing: the price table
itself (O9) and the set of channels required to settle through $\mathsf{sink}$ (O8).

% ---------------------------------------------------------------------------
\subsection{Settlement shapes not taken}
\label{sec:pnr:alternatives}

Two other shapes were considered against the same constraints --- the coinbase must not grow, not
every wallet exposes multiple recipients, adoption should be fast --- and both were rejected. They
are recorded because a reader will assume one of them was the easy option
\srcdoc{plan/conflicts.md C36}.

\paragraph{Not a blockchain fee.} Carrying the operator share as transaction fee is the shape that
looks free, and it is not, for four independent reasons. \emph{First}, fees accrue to the block
producer, so routing them into a pool instead would be a consensus change anyway --- the same
implementation cost as a sink, with worse properties, since a producer that includes a payment would
capture part of it in the same block, which is exactly what Theorem~\ref{thm:pnr:timing} exists to
remove. \emph{Second}, the paid amount would become
$\sum\text{inputs} - \sum\text{outputs}$, a quantity the sending wallet computes, so a change-output
calculation would silently determine what the tenant paid --- and a payment amount that no party
states explicitly is not a payment a tenant can audit. \emph{Third}, wallets clamp, warn on, or
refuse anomalously large fees, and no fee-estimation interface is designed for a fee that is
\SI{20}{\percent} of a deployment price, so the shape would fight the tooling on every payment.
\emph{Fourth}, an ordinary transaction carrying a large fee is indistinguishable from an application
payment except by its memo, so the memo would become load-bearing for \emph{consensus} rather than
for attribution --- a strictly stronger requirement than PS2 makes, and one that puts a
\FluxOS-defined field inside a money rule.

\paragraph{Not a new transaction type, at least not first.} A typed transaction is the most rigorous
shape, and \Flux has the in-house precedent: transaction versions 5 and 6 and the \FluxNode
start/confirm types \shipped \srccode{fluxd/\allowbreak{}src/\allowbreak{}fluxnode/\allowbreak{}activefluxnode.cpp}{409--427}. It was not
chosen as the primary path for two reasons. It would force every payer to link a \Flux-specific
library in order to pay, which is precisely the friction \S18 exists to remove and which no ordinary
wallet or agent runtime would carry; and it would have to be rolled out across every wallet, hosted
tool and agent before any payment could settle at all, making adoption serial rather than immediate.
It would remain available as a later addition carrying richer fields, and it would coexist with the
address-and-memo path without invalidating it \planned \srcdoc{plan/conflicts.md C36}.

% ---------------------------------------------------------------------------
\subsection{The pool as consensus state}
\label{sec:pnr:pool}

\begin{definition}[Fee pool and per-block transition, \planned]
\label{def:pnr:pool}
The \emph{fee pool} is a single consensus-state integer $\pool_\hgt \in \mathbb{Z}_{\ge 0}$ in
$\sat$, with $\pool_{\hpa - 1} \defeq 0$. Its value would be \emph{committed in the coinbase} of
block $\hgt$ and recomputed by every validator from the parent value and the block body (rule PC
below). Write $\Delta_\hgt \defeq \pool_\hgt - \pool_{\hgt-1}$ for the per-block delta, which an
implementation may cache against the block index but which carries no authority: the pool is a pure
function of the chain. Fix the \emph{drip constant} $\drip \in \mathbb{Z}_{\ge 1}$, denominated in
blocks and settled at $\drip = 86{,}400$ (\S\ref{sec:pnr:residual}), and write $\mathcal{O}_\hgt$
for the set of outputs paying $\mathsf{sink}$ in block $\hgt$. For every block $\hgt \ge \hpa$:
\begin{align}
  D_\hgt      &\defeq \floorfrac{\pool_{\hgt-1}}{\drip}
                 &&\text{(drip, taken from the \emph{parent} state)}, \label{eq:pnr:drip}\\
  D_{\hgt,t}  &\defeq \floorfrac{2\tbase{t}\,D_\hgt}{27}, \quad t \in \tiers
                 &&\text{(tier partition, over the integers $2\tbase{t} \in \{2,7,18\}$)},
                 \label{eq:pnr:tier}\\
  \epsilon_\hgt &\defeq D_\hgt - \textstyle\sum_{t \in \tiers} D_{\hgt,t} \in \{0,1,2\}\sat
                 &&\text{(partition residue, retained in the pool)}, \label{eq:pnr:eps}\\
  F_\hgt      &\defeq \textstyle\sum_{o \in \mathcal{O}_\hgt} A_\Phi(o)
                 &&\text{(inflow, over the $\mathsf{sink}$-paying outputs of block $\hgt$)},
                 \label{eq:pnr:inflow}\\
  G_\hgt      &\defeq \textstyle\sum_{o \in \mathcal{O}_\hgt} A_F(o)
                 &&\text{(Foundation share of the same receipts)}, \label{eq:pnr:found}\\
  \pool_\hgt  &\defeq \pool_{\hgt-1} - D_\hgt + \epsilon_\hgt + U_\hgt + F_\hgt
                 &&\text{(update)}. \label{eq:pnr:update}
\end{align}
For each tier $t$ with $\payees_\hgt^{t} \ne \varnothing$, the existing tier output to
$\payees_\hgt^{t}$ carries value $\subsidy_t(\hgt) + D_{\hgt,t}$; the existing dev-fund output
carries $\devfund(\hgt) + G_\hgt$. One further rule closes the state:
\begin{description}
  \item[PC (pool commitment).] The coinbase of every block $\hgt \ge \hpa$ carries $\pool_\hgt$ as
    an \num{8}-byte little-endian integer in $\sat$, in the coinbase's data-bearing field --- its
    input signature script, which holds data rather than value, so the coinbase gains no output
    (Property~\ref{prop:pnr:invariants}, I4). A validator recomputes $\pool_\hgt$ by
    \eqref{eq:pnr:drip}--\eqref{eq:pnr:update} and rejects the block if the committed value differs.
\end{description}
$U_\hgt$ is the sum of $D_{\hgt,t}$ over tiers with $\payees_\hgt^{t} = \varnothing$, retained in the
pool. The tier weights are $\tierratio{t} = \tbase{t}/\sum_{u}\tbase{u}$ with
$(\tbase{\tC}, \tbase{\tN}, \tbase{\tS}) = (1,\ 3.5,\ 9)\flux$
\srccode{fluxd/\allowbreak{}src/\allowbreak{}main.cpp}{2886--2939}, giving
$(\tierratio{\tC}, \tierratio{\tN}, \tierratio{\tS}) = (2/27,\ 7/27,\ 18/27)
 = (\SI{7.41}{\percent},\ \SI{25.93}{\percent},\ \SI{66.67}{\percent})$.
\end{definition}

\settlednote{Stays in $\pool$ (\Cref{tab:pnr-recommended}). O5 --- destination of $D_{\hgt,t}$ when tier $t$ has no confirmed node. Domain:
$\{\text{stays in pool}$, $\text{dev fund}$, $\text{redistributed to the other tiers}\}$.
Definition~\ref{def:pnr:pool} takes the conservation-trivial default (pool) so that the invariants
of \S\ref{sec:pnr:invariants} hold as written; the dev-fund option would mirror the subsidy fallback
reported for \PoN, and the redistribution option rewards concentration.}

\settlednote{Stays in $\pool$ (\Cref{tab:pnr-recommended}). O6 --- destination of the partition residue $\epsilon_\hgt \le 2\sat$ per block. Domain:
$\{\text{pool}, \text{dev fund}\}$. Worth at most $0.021\flux$ per year either way; it must be fixed
explicitly because an implementation difference here is a chain split
(\S\ref{sec:pnr:invariants}, I5).}

\settlednote{Drip from $\pool_{\hgt-1}$, the default kept (\Cref{tab:pnr-recommended}). O10 --- drip ordering. Domain: $\{\text{drip from } \pool_{\hgt-1},\ \text{drip after
applying } F_\hgt\}$. Definition~\ref{def:pnr:pool} drips from the parent state, which buys two
properties used later: the coinbase is checkable before the block body is executed, and a block
producer cannot influence its own drip. Dripping after inflow would remove one block of lag and
would reintroduce same-block recapture of $\tierratio{t}/\drip$.}

\begin{remark}[O21, settled: the division is applied on arrival]
\label{rem:pnr:o21}
The settlement decision created a question this specification had not needed to ask --- at which
height the Foundation share $G_\hgt$ is computed --- and it is answered here: \textbf{on arrival},
at the $\nodeshare$ in force at the receipt height, as PS4 states, and paid in the same block
\srcintent{evidence/design/08-answers-round-2.md, round 18}.

The alternative considered was to let the pool receive the entire receipt and apply the division to
the drip at payout, adding $\lfloor (100 - r(\hgt))\,D_\hgt/100 \rfloor$ to the existing dev-fund
output. It is attractive on its face: it needs \emph{one} consensus integer rather than two, and
every coinbase output would remain checkable from the parent state alone, because $D_\hgt$ depends
only on $\pool_{\hgt-1}$.

It is rejected for a reason that is easy to miss. $\nodeshare$ is not constant --- it ramps by
$\rampstep$ at each subsidy reduction --- and $\pool$ is a single integer that cannot carry the rate
that applied when each receipt arrived. Applying the division at payout therefore re-prices the
entire accumulated balance at every ramp step: of order $\rampstep\,\pool$, once per reduction,
moved from the Foundation to operators on revenue earned under the earlier share. The ramp is a
schedule for dividing revenue \emph{as it is earned}; deferring the division to payout silently
converts it into a schedule over \emph{distributions}. That is a different policy, and not the one
the ramp was specified to express.

What is given up is less than it appears. Parent-state checkability of the coinbase serves a
verifier that does not hold the block body, and \Cref{sec:bridge} establishes that no such verifier
is constructible for \Flux. A dev-fund output that depends on the receipts in its own block is in
any case the ordinary arrangement wherever transaction fees are paid.
\end{remark}

\begin{construction}[\PNR block connect, \planned]
\label{con:pnr:connect}
\textsc{Inputs}: parent pool $\pool_{\hgt-1}$, read from the parent block's coinbase commitment; the
block at height $\hgt$ with its transaction list; the \PoN payee tuple
$\payees_\hgt = (\payees_\hgt^{\tC}, \payees_\hgt^{\tN}, \payees_\hgt^{\tS})$ returned by the
existing queue; the constants $\drip$, $\{\tbase{t}\}$, $\mathsf{sink}$, $\nodeshare(\hgt)$.\\
\textsc{Invariant} (maintained across every call):
$\pool_\hgt = \sum_{j \le \hgt} F_j - \sum_{j \le \hgt}\big(D_j - \epsilon_j - U_j\big) \ge 0$.\\
\textsc{Failure mode}: any step that does not hold makes the block invalid; there is no partial
application. A divergent implementation is detected at $\hgt$ by step 6, not one block later. On
\texttt{Disconnect\allowbreak{}Block} the cached $\Delta_\hgt$ is subtracted, restoring $\pool_{\hgt-1}$ exactly;
if the cache is absent or suspect, $\pool_{\hgt-1}$ is read from the parent coinbase or recomputed
from $\hpa$, and there is no repair path to invoke because there is nothing that can diverge from
the chain.
\begin{enumerate}
  \item Read $\pool_{\hgt-1}$. If $\hgt < \hpa$, set $\Delta_\hgt \defeq 0$ and stop.
  \item Compute $D_\hgt$ by \eqref{eq:pnr:drip}, then $D_{\hgt,t}$ and $\epsilon_\hgt$ by
    \eqref{eq:pnr:tier}--\eqref{eq:pnr:eps}. All three depend only on the parent state and on
    constants.
  \item Scan the block's outputs for payments to $\mathsf{sink}$. For each such output of value $v$,
    compute $A_\Phi$ and $A_F$ by PS4 of Definition~\ref{def:pnr:tx} at this height; accumulate
    $F_\hgt$ and $G_\hgt$. Reject the block if any transaction in it spends an output paying
    $\mathsf{sink}$ (PS3).
  \item For each tier $t$: if $\payees_\hgt^{t} \ne \varnothing$, require the coinbase tier output
    to equal exactly $\subsidy_t(\hgt) + D_{\hgt,t}$; otherwise add $D_{\hgt,t}$ to $U_\hgt$.
    Require the dev-fund output to equal $\devfund(\hgt) + G_\hgt$.
  \item Set $\Delta_\hgt \defeq F_\hgt - D_\hgt + \epsilon_\hgt + U_\hgt$ and
    $\pool_\hgt \defeq \pool_{\hgt-1} + \Delta_\hgt$; cache $\Delta_\hgt$ against the block index.
  \item Require the coinbase commitment PC to equal $\pool_\hgt$; reject the block on mismatch.
  \item If the block is an emergency block (Assumption~\ref{as:pnr} A3), step 2 is skipped and
    $D_\hgt \defeq 0$, so $D_{\hgt,t} = \epsilon_\hgt = 0$ and the tier clause of step 4 reduces to
    today's rule --- see Property~\ref{prop:pnr:exported}.
\end{enumerate}
\end{construction}

\paragraph{Where the pool lives, and the two places it must not.} O2 asked whether $\pool_\hgt$ is
held as a per-index delta, a coinbase field or a header field; the answer is the coinbase field, PC
above, with every validator recomputing rather than trusting it \srcdoc{plan/conflicts.md C37}. Three
properties decide it. \emph{Reorg-exactness would come free}, because the pool is a pure function of
the chain: a reorg recomputes it, with no rollback journal and no possibility of the pool diverging
from the UTXO set. \emph{No new database would be needed}, because validators already parse the
coinbase for the dev-fund minimum and the three tier payouts \shipped
\srccode{fluxd/\allowbreak{}src/\allowbreak{}main.cpp}{1231--1249}, \srccode{fluxd/\allowbreak{}src/\allowbreak{}fluxnode/\allowbreak{}fluxnode.cpp}{837--875}, so the
parsing and enforcement path exists and is exercised every block. And \emph{anyone could check it
from chain data alone}, which is what \S5's supply-transparency argument and
Property~\ref{prop:pnr:supply} both require.

A \emph{separate LevelDB store} is rejected on this network's own precedent. The \FluxNode registry
lives in its own store and needed a dedicated crash-recovery marker, precisely because the coins
database and the node database can diverge after an unclean shutdown, together with a repair RPC to
rebuild it \shipped \srccode{fluxd/\allowbreak{}src/\allowbreak{}fluxnode/\allowbreak{}fluxnodecachedb.h}{22--24},
\srccode{fluxd/\allowbreak{}src/\allowbreak{}rpc/\allowbreak{}fluxnode.cpp}{2676}. That pattern is tolerable for a registry that can be
rebuilt from chain data. It is not tolerable for money, where a divergence between the store and the
chain is a wrong payout rather than a stale index.

A \emph{header field} is rejected for a different reason. It is the most rigorous option, but a
header change is the widest-blast-radius fork available to this chain, and it would further
complicate the light-client position of \S22 --- which is already foreclosed by the absence of a
node-set commitment, by fixed $2^{40}$ chainwork (Assumption~\ref{as:pnr} A2) and by the emergency
path (A3) (C26). One circumstance would reverse this: if a node-set commitment is ever added to the
header to make light clients constructible (\S8, \S22), that is the moment to reconsider carrying
the pool commitment alongside it, since the fork is being taken anyway.

\begin{proposition}[Closed form, steady state and half-life, \planned]
\label{prop:pnr:closed}
Relax the floors in \eqref{eq:pnr:drip}--\eqref{eq:pnr:eps} (they retain at most $3\sat$ per block in
the pool and therefore only \emph{understate} the pool). Write $q \defeq 1 - 1/\drip \in (0,1)$ for
$\drip \ge 2$. Then for all $\hgt \ge \hgt_0 \ge \hpa$,
\[
  \pool_\hgt \;=\; q^{\,\hgt - \hgt_0}\pool_{\hgt_0}
    \;+\; \sum_{j = \hgt_0+1}^{\hgt} q^{\,\hgt - j} F_j .
\]
Under constant inflow $F_j \equiv F$ this gives
$\pool_\hgt = q^{\hgt}\pool_0 + F\drip\,(1 - q^{\hgt})$, hence
\[
  \pool^{*} \;\defeq\; \lim_{\hgt \to \infty}\pool_\hgt \;=\; F\,\drip ,
  \qquad
  D^{*} \;=\; \pool^{*}/\drip \;=\; F .
\]
The half-life of a perturbation --- equivalently, of the pool under zero inflow --- is
\[
  \hgt_{1/2} \;=\; \frac{\ln 2}{-\ln(1 - 1/\drip)} \;=\; \drip \ln 2 + \tfrac12 \ln 2 + O(1/\drip).
\]
A single contribution $A_\Phi$ made in block $\hgt_0$ is paid out as $A_\Phi\,q^{j}/\drip$ in block
$\hgt_0 + 1 + j$ for $j \ge 0$, and $\sum_{j \ge 0} A_\Phi q^{j}/\drip = A_\Phi$.
\end{proposition}

\begin{proof}
With the floors relaxed, \eqref{eq:pnr:update} reads $\pool_\hgt = q\,\pool_{\hgt-1} + F_\hgt$;
unrolling from $\hgt_0$ gives the stated sum. For $F_j \equiv F$ the sum is geometric,
$F\sum_{i=0}^{\hgt - \hgt_0 - 1} q^{i} = F\drip(1 - q^{\hgt - \hgt_0})$ since
$1 - q = 1/\drip$. Letting $\hgt \to \infty$ with $|q| < 1$ gives $\pool^{*} = F\drip$, and
$D^{*} = \pool^{*}/\drip = F$. For the half-life, solve $q^{n} = 1/2$ for $n$, giving
$n = \ln 2 / (-\ln q)$; expanding $-\ln(1 - 1/\drip) = 1/\drip + 1/(2\drip^{2}) + O(\drip^{-3})$
gives $n = \drip\ln 2 \,(1 + 1/(2\drip) + O(\drip^{-2}))$, i.e.\ the stated expansion. The impulse
response is the $F_j$ term of the closed form with a single non-zero $F_{\hgt_0} = A_\Phi$, and its
sum over $j$ is $A_\Phi \cdot (1/\drip)\cdot (1-q)^{-1} = A_\Phi$.
\end{proof}

\begin{remark}[The pool is a delay line, not a treasury]
\label{rem:pnr:delayline}
Proposition~\ref{prop:pnr:closed} says that at steady state the drip equals the inflow: everything
paid in is paid out, with lag. Nothing accumulates permanently, nothing is burned, and there is no
balance that anyone could later choose how to spend --- the value held at $\mathsf{sink}$ is not a
treasury either, because PS3 makes it unspendable by every party including the Foundation, and its
only exit is the coinbase schedule. The pool is a delay line with time constant
$\drip$ blocks. This framing matters at scale: at parity-scale revenue
(\S\ref{sec:pnr:crossover}, $31.9$M$\flux$ per year at $\nodeshare = 0.5$) the accumulator would
hold about $1.31$M$\flux$ at the recommended $\drip$, and a reader who mistakes that number for a
treasury will draw exactly the wrong conclusion about who controls it. Nobody controls it; it is
$\SI{30}{\day}$ of pipeline.
\end{remark}

At $\drip = 86{,}400$ blocks and $\spacingpon = \SI{30}{\second}$
\srccode{fluxd/\allowbreak{}src/\allowbreak{}chainparams.cpp}{105}, the fill from zero would be
$D_\hgt/F = 1 - q^{\hgt - \hpa}$: \SI{3.3}{\percent} after one day, \SI{20.8}{\percent} after a
week, \SI{63.2}{\percent} after $\drip$ blocks, \SI{95.0}{\percent} after \SI{90}{\day}
\srcdoc{plan/mech-pnr.md \S1.5}. \PNR would therefore ramp over roughly a quarter rather than switch
on, and the paper should not describe activation as a switch.

\settlednote{No seed: $\pool_{\hpa-1}=0$ (\Cref{tab:pnr-recommended}). O11 --- initial seed $\pool_{\hpa-1}$. Domain: $\{0,\ \text{a Foundation-funded seed}\}$. A
seed removes the ninety-day ramp; it is also a one-off Foundation transfer to operators and would
have to be labelled as one.}

% ---------------------------------------------------------------------------
\subsection{Fairness over a rotation}
\label{sec:pnr:fairness}

\begin{proposition}[Rotation, \shipped]
\label{prop:pnr:rotation}
Under the queue rule of Assumption~\ref{as:pnr} A1, with a static confirmed set of $\ncount{t}$
nodes in tier $t$, the payee sequence of tier $t$ is periodic with period $\rot{t} = \ncount{t}$ and
every node in the tier is paid exactly once per period.
\end{proposition}

\begin{proof}
Order the tier's nodes by their sort key (last-paid height, or confirmation height for a node never
paid). The front node is paid at block $\hgt$ and its key becomes $\hgt$, which strictly exceeds
every other key, since every other node was either paid at some height $< \hgt$ or confirmed at some
height $< \hgt$. It therefore moves to the back, and the relative order of the remaining
$\ncount{t} - 1$ nodes is unchanged. Each block thus rotates the sequence by exactly one position;
after $\ncount{t}$ blocks the sequence is the original. Hence the period is $\ncount{t}$ and each
node is paid exactly once within it.
\end{proof}

\begin{remark}[Churn, and what the chain is actually doing]
\label{rem:pnr:churn}
A node joining is keyed by its confirmation height and enters at the back; a node leaving --- by
spending its collateral, or by failing to re-confirm within $640$ blocks --- shortens the cycle
\shipped \srccode{fluxd/\allowbreak{}src/\allowbreak{}fluxnode/\allowbreak{}fluxnode.h}{31--34}. In a window of $\ncount{t}$ blocks a node
is therefore paid at most $1 + (\text{departures})$ times. Proposition~\ref{prop:pnr:rotation} is
not a modelling assumption: at chain tip $2{,}917{,}115$ the maximum observed
$\text{tip} - \text{last-paid height}$ was $3{,}202$ against $3{,}247$ Cumulus nodes, $1{,}597$
against $1{,}615$ Nimbus, and $1{,}631$ against $1{,}627$ Stratus
\srcmeasured{api.runonflux.io/\allowbreak{}daemon/\allowbreak{}listzelnodes}{2026-09-03}. Cumulus and Nimbus sit inside one
rotation; Stratus exceeds it by four blocks, consistent with four joins during that cycle at the
measured median tenure of $\SI{67}{\day}$. Rotation lengths at $\SI{30}{\second}$ are
$\SI{27.06}{\hour}$, $\SI{13.46}{\hour}$ and $\SI{13.56}{\hour}$ respectively.
\end{remark}

\begin{proposition}[Equal income within a tier, \planned]
\label{prop:pnr:within}
Assume Proposition~\ref{prop:pnr:rotation} and zero churn. Over $M$ full rotations every node in
tier $t$ would receive exactly $M$ \PNR payments, and for any two nodes $i, j$ in the tier with
cumulative \PNR incomes $I_i, I_j$,
\[
  \frac{\big|I_i - I_j\big|}{\min(I_i, I_j)} \;\le\; s_t
  \;\defeq\; \max_{\text{rotation}} \frac{\max_{\hgt} D_\hgt}{\min_{\hgt} D_\hgt} - 1 ,
\]
where the maximum and minimum are taken over the blocks of one rotation. The mixing bound is one
rotation, i.e.\ $\rot{t}$ blocks.
\end{proposition}

\begin{proof}
By Proposition~\ref{prop:pnr:rotation} each rotation gives each node of tier $t$ exactly one payment,
of size $\tierratio{t} D_{\hgt}$ for the block $\hgt$ at which it is the payee. Within one rotation
any two such payments differ by at most a factor $1 + s_t$ by the definition of $s_t$. Summing $M$
rotations, $I_i = \sum_{k=1}^{M}\tierratio{t}D_{\hgt_k^{i}}$ and likewise for $j$, and a
term-by-term comparison preserves the factor. Dividing by $\min(I_i, I_j)$ gives the claim.
\end{proof}

\begin{remark}[The bound does not vanish with averaging]
\label{rem:pnr:phase}
A node's phase within the rotation is fixed --- it is paid at the same position every cycle --- so an
inflow pattern periodic with the rotation could favour one phase systematically. Averaging over many
rotations therefore does \emph{not} drive the difference to zero; it remains bounded by $s_t$ and no
better. This is why $s_t$ must be small in absolute terms, which is the subject of
\S\ref{sec:pnr:residual}, rather than small on average.
\end{remark}

\begin{proposition}[Cross-tier income is issuance-shaped, not collateral-proportional, \planned]
\label{prop:pnr:across}
Per-block-averaged \PNR income per node in tier $t$ is $\tierratio{t}\bar{D}/\ncount{t}$, where
$\bar{D}$ is the mean drip. Hence for tiers $t, t'$,
\[
  \frac{I_t}{I_{t'}} \;=\; \frac{\tbase{t}}{\tbase{t'}}\cdot\frac{\ncount{t'}}{\ncount{t}} ,
  \qquad\text{and}\qquad
  \frac{\text{\PNR income}}{\text{\PoN issuance income}}
  \;=\; \frac{\bar{D}}{\tfrac{13.5}{14}\,\subsidypon}
\]
for \emph{every} node, the same constant across all tiers.
\end{proposition}

\begin{proof}
By Proposition~\ref{prop:pnr:rotation} each node of tier $t$ is paid once every $\ncount{t}$ blocks,
and its payment is $\tierratio{t}$ of that block's drip; averaging over blocks gives
$\tierratio{t}\bar D/\ncount{t}$ per node per block. Taking the ratio for two tiers, $\bar D$
cancels and $\tierratio{t}/\tierratio{t'} = \tbase{t}/\tbase{t'}$ by
Definition~\ref{def:pnr:pool}. For the second identity, \PoN issuance income per node in tier $t$
is likewise $\tierratio{t}\cdot\tfrac{13.5}{14}\subsidypon/\ncount{t}$ per block, because the drip
and the subsidy are paid to the same payees in the same tier proportions and the dev-fund remainder
is $0.5/14$ of the subsidy \srccode{fluxd/\allowbreak{}src/\allowbreak{}main.cpp}{2877--2884}; the ratio of the two
expressions is independent of $t$.
\end{proof}

\begin{remark}[The correction, stated without softening]
\label{rem:pnr:c24}
Substituting the measured fleet into Proposition~\ref{prop:pnr:across}:
Nimbus-to-Cumulus is $3.5 \times 3{,}246/1{,}615 = 7.04$, and Stratus-to-Cumulus is
$9 \times 3{,}246/1{,}627 = 17.96$, i.e.\ \textbf{a Stratus node would earn $18.0\times$ a Cumulus
node while posting $40\times$ the collateral} ($\coll{\tC} = 1{,}000\flux$,
$\coll{\tN} = 12{,}500\flux$, $\coll{\tS} = 40{,}000\flux$
\srccode{fluxd/\allowbreak{}src/\allowbreak{}fluxnode/\allowbreak{}fluxnode.h}{57--63}). Return on collateral relative to Cumulus is
$7.04/12.5 = 0.563$ for Nimbus and $17.96/40 = 0.449$ for Stratus: \textbf{per \flux of collateral,
Cumulus is the better-compensated tier, by a factor of $2.23$ over Stratus.} The design intake's
phrase ``proportional to collateral across tiers'' is therefore incorrect and is not used anywhere
in this paper \srcintent{evidence/design/02-pnr.md} (C24). What \PNR income \emph{is} exactly
proportional to is \PoN issuance income --- it is issuance-shaped --- so the tilt toward Cumulus on a
return-on-collateral basis is a pre-existing property of \PoN that \PNR would inherit and reinforce.
\end{remark}

% ---------------------------------------------------------------------------
\subsection{Residual unfairness, quantified}
\label{sec:pnr:residual}

Proposition~\ref{prop:pnr:within} bounds within-tier inequality by $s_t$ but does not make it zero,
because the pool decays between the first and last payee of a rotation. Two effects set $s_t$.

\begin{proposition}[Intra-rotation spread, \planned]
\label{prop:pnr:spread}
Within a rotation of length $\rot{}$ under zero inflow, the node paid at position $i \in
\{0,\dots,\rot{}-1\}$ receives $\tierratio{t}\pool_{\hgt_0}q^{i}/\drip$, so first-versus-last payee
gives
\[
  s^{\downarrow} \;=\; q^{-(\rot{}-1)} - 1 \;\approx\; e^{\rot{}/\drip} - 1 .
\]
A single contribution $A_\Phi$ arriving mid-rotation raises the pool multiplicatively, so the node
paid just after it receives $(1 + A_\Phi/\pool)$ times the node paid just before; at steady state
$\pool^{*} = F\drip$ this upside jump is $\nodeshare\price/(F\drip)$. Combining both sides, with all
of a rotation's inflow arriving as a lump at its end,
\[
  s \;\le\; e^{\rot{}/\drip}\big(1 + \rot{}/\drip\big) - 1 \;\approx\; 2\rot{}/\drip .
\]
\end{proposition}

\begin{proof}[Proof sketch]
The decay term is Proposition~\ref{prop:pnr:closed} with $F \equiv 0$: $\pool$ is multiplied by $q$
each block, so payments at positions $0$ and $\rot{}-1$ stand in ratio $q^{-(\rot{}-1)}$, and
$q^{-n} = e^{-n\ln(1-1/\drip)} = e^{n/\drip + O(n/\drip^{2})}$. The upside term is immediate from
$D_\hgt = \pool_{\hgt-1}/\drip$ being linear in $\pool$. The combined bound multiplies the worst
decay factor by the worst single-lump jump $1 + \rot{} F/(F\drip)$ and is loose by construction: it
assumes the entire rotation's inflow arrives in one block at the least favourable moment. A tight
bound would require a distributional assumption on $F_\hgt$ that this specification does not make.
\end{proof}

Evaluated against the longest rotation, Cumulus at $\rot{\tC} = 3{,}246$ blocks
\srcdoc{paper/data/pnr\_drip.dat}:

\begin{table}[H]
\centering
\small
\begin{tabular}{rrrrrr}
\toprule
$\drip$ (blocks) & $\hgt_{1/2}$ (days) & pool decay per & $s^{\downarrow}$ (first vs.\ & loose steady- & upside jump of one\\
 & & Cumulus rotation & last, zero inflow) & state bound & $1{,}000\flux$ payment at $V_0$\\
\midrule
$3{,}246$   & $0.78$  & \SI{63.22}{\percent} & \SI{171.79}{\percent} & \SI{443.7}{\percent} & \SI{22.6}{\percent}\\
$10{,}000$  & $2.41$  & \SI{27.72}{\percent} & \SI{38.34}{\percent}  & \SI{83.3}{\percent}  & \SI{7.3}{\percent}\\
$32{,}460$  & $7.81$  & \SI{9.52}{\percent}  & \SI{10.51}{\percent}  & \SI{21.6}{\percent}  & \SI{2.3}{\percent}\\
$64{,}920$  & $15.62$ & \SI{4.88}{\percent}  & \SI{5.13}{\percent}   & \SI{10.4}{\percent}  & \SI{1.1}{\percent}\\
$\mathbf{86{,}400}$ & $\mathbf{20.79}$ & \textbf{\SI{3.69}{\percent}} & \textbf{\SI{3.83}{\percent}} & \SI{7.7}{\percent} & \textbf{\SI{0.85}{\percent}}\\
\bottomrule
\end{tabular}
\caption{Residual within-tier unfairness as a function of the drip constant $\drip$, against the
measured Cumulus rotation $\rot{\tC} = 3{,}246$ blocks. Steady-state pool at $V_0$ and
$\nodeshare = 0.20$: $\pool^{*} = F_0\drip$ with $F_0 = 0.27237\flux$ per block, giving
$884 / 2{,}724 / 8{,}841 / 17{,}683 / 23{,}533\flux$ down the column. Computed by
\texttt{scripts/\allowbreak{}07-pnr-settlement.\allowbreak{}py}, 2026-09-03.}
\label{tab:pnr:drip}
\end{table}

At activation-scale revenue the \emph{upside} term, not the decay term, would dominate the realised
spread: with $\pool^{*} \approx 23{,}533\flux$ at $\drip = 86{,}400$, a $100\flux$ payment moves the
drip \SI{0.085}{\percent}, a $1{,}000\flux$ payment \SI{0.85}{\percent}, a $10{,}000\flux$ payment
\SI{8.5}{\percent}, and a $100{,}000\flux$ payment \SI{85}{\percent} --- and deployments at the
upper end are constructible, since $\ninst \le 100$ and $\dur \le 1{,}056{,}000$ blocks
\srcmeasured{api.runonflux.io/\allowbreak{}apps/\allowbreak{}deploymentinformation}{2026-09-03}. The decay term is what a
larger $\drip$ controls; the upside term shrinks only as revenue and therefore $\pool^{*}$ grow. At
parity-scale revenue the same $10{,}000\flux$ payment would move the pool \SI{0.4}{\percent}. Both
statements belong together, and neither should be quoted without the other.

Fleet growth at fixed $\drip = 86{,}400$ raises $s^{\downarrow}$: \SI{3.83}{\percent} at $3{,}246$
Cumulus nodes, \SI{5.96}{\percent} at $5{,}000$, \SI{7.80}{\percent} at $2\times$,
\SI{11.93}{\percent} at $3\times$, \SI{16.21}{\percent} at $4\times$, and \SI{26.05}{\percent} at
$20{,}000$ \srcdoc{scripts/07-pnr-settlement.py}.

\paragraph{The value, and what it costs.} $\drip = 86{,}400$ blocks --- exactly $\SI{30}{\day}$ at
$\SI{30}{\second}$ spacing --- giving a decay-side spread of \SI{3.83}{\percent}, a
$\SI{20.79}{\day}$ half-life, and the recapture bound proved in \S\ref{sec:pnr:timing}, with
headroom for fleet growth. The costs are stated rather than netted: demand changes would reach
operators with a thirty-day time constant; \PNR would ramp from zero over about ninety days at
activation unless seeded; and the accumulator would hold roughly $1.3$M$\flux$ at parity-scale
revenue, which is a delay line and not a treasury (Remark~\ref{rem:pnr:delayline}). The value
$\drip = \rot{\tC}$ that the intake first suggested is rejected on the \SI{171.79}{\percent} figure
in Table~\ref{tab:pnr:drip}. $\drip$ is not to be adaptive to node count: node count is itself
gameable by registration, and a fixed constant with margin dominates
\srcintent{evidence/design/02-pnr.md}.

\paragraph{Ratified.} $\drip = 86{,}400$ blocks is settled and O1 is closed
\srcintent{evidence/design/08-answers-round-2.md}. The trade-off it resolves does not disappear with
the decision and is restated once so that a later revision starts from it: spread
$\approx \rot{\tC}/\drip$ and targeted recapture $\approx \tierratio{t}/\drip$ both fall with
$\drip$, while lag, activation ramp and steady-state pool size $F\drip$ all rise with it, over the
domain $[10^{4},\ 2\times10^{5}]$ blocks in which the constant was chosen. What remains open is not
the constant but the guard on it under fleet growth (O12), and the constant is fixed by consensus,
so revisiting it would be a network upgrade rather than a configuration change.

\settlednote{A documented review trigger at $2\times$, not an automatic rule (\Cref{tab:pnr-recommended}). O12 --- fleet-growth guard on $\drip$. Domain: $\{\text{none},\ \text{a documented
hard-fork review trigger}\}$. At $4\times$ the current Cumulus count the decay-side spread reaches
\SI{16}{\percent}. An automatic adaptive rule is excluded (it is a manipulation surface); a
documented trigger is not.}

% ---------------------------------------------------------------------------
\subsection{The timing attack, and why the drip defeats it}
\label{sec:pnr:timing}

\begin{definition}[Payout-timing attack, \planned]
\label{def:pnr:attack}
Let $\adv$ be an operator controlling $m_t$ confirmed nodes in tier $t$, with a capital budget
$\advbudget$ sufficient to have registered them, and let $\adv$ also be a tenant --- it deploys its
own applications and pays for them. By Assumption~\ref{as:pnr} A1 the payee sequence is a public
deterministic function of chain state, so $\adv$ knows every future height $\hgt$ at which one of its
nodes satisfies $\payees_\hgt^{t} \in \adv$. The \emph{payout-timing attack} is: submit a payment of
size $\price$ so that it is included in block $\hgt - 1$, where $\hgt$ is such an aligned height, and
collect whatever share of that payment the settlement rule returns at $\hgt$. The attack costs
nothing --- $\adv$ times a payment it would have made anyway --- and requires waiting at most one
rotation, i.e.\ at most $\SI{27.06}{\hour}$, for an aligned height.
\end{definition}

\begin{proposition}[Recapture under next-block payout, \planned]
\label{prop:pnr:nextblock}
Under the intake's first-cut settlement rule --- fees collected in block $\hgt-1$ distributed in
block $\hgt$ to $\payees_\hgt$ by tier ratio --- the attack of Definition~\ref{def:pnr:attack}
recaptures
\[
  \mathrm{recap}_{\mathrm{next}}
  \;=\; \nodeshare\,\price \sum_{t \in \tiers} \tierratio{t}\,
        \indic{\payees_\hgt^{t} \in \adv}.
\]
\end{proposition}

\begin{proof}
Immediate from the rule: the operator share $\nodeshare\price$ of the block-$(\hgt-1)$ payments is
split across the three tier payees of block $\hgt$ in proportions $\tierratio{t}$, and $\adv$
collects the terms whose payee it controls.
\end{proof}

\begin{table}[H]
\centering
\small
\begin{tabular}{lrrr}
\toprule
$\adv$ controls & of the \emph{operator share} & of the payment, $\nodeshare = 0.20$ & of the payment, $\nodeshare = 0.50$\\
\midrule
one Cumulus node          & \SI{7.41}{\percent}  & \SI{1.48}{\percent}  & \SI{3.70}{\percent}\\
one Nimbus node           & \SI{25.93}{\percent} & \SI{5.19}{\percent}  & \SI{12.96}{\percent}\\
one Stratus node          & \textbf{\SI{66.67}{\percent}} & \textbf{\SI{13.33}{\percent}} & \SI{33.33}{\percent}\\
one node of each tier, aligned & \textbf{\SI{100}{\percent}} & \textbf{\SI{20}{\percent}} & \SI{50}{\percent}\\
\bottomrule
\end{tabular}
\caption{Targeted-block recapture under next-block payout. \textbf{The intake's headline figures
of \SI{66.7}{\percent} and \SI{100}{\percent} are fractions of the \emph{operator share}, not of the
payment}; as fractions of what the tenant pays they are \SI{13.3}{\percent} and \SI{20}{\percent} at
activation, rising to \SI{33.3}{\percent} and \SI{50}{\percent} at the cap
\srcintent{evidence/design/02-pnr.md}.}
\label{tab:pnr:nextblock}
\end{table}

With corrected magnitudes the intake's conclusion stands: every operator holding a Stratus node
would obtain a \SIrange{13}{33}{\percent} discount on its own hosting, and an operator holding one
node of each tier would host free at activation. The operator share would become a rebate to
operator-tenants rather than income. Next-block payout must not ship.

\begin{theorem}[Timing resistance of the drip, \planned]
\label{thm:pnr:timing}
Under Definition~\ref{def:pnr:pool} the attack of Definition~\ref{def:pnr:attack} recaptures, in the
targeted block $\hgt$, exactly
\[
  \frac{\mathrm{recap}_{\mathrm{drip}}(\hgt)}{A_\Phi} \;=\; \frac{\tierratio{t}}{\drip},
  \qquad A_\Phi = \nodeshare\,\price .
\]
For a single Stratus node at $\drip = 86{,}400$ this is
$\tierratio{\tS}/\drip = (2/3)/86{,}400 = 7.72 \times 10^{-6}$, i.e.\
\textbf{\SI{0.00077}{\percent} of the operator share} and $1.54\times10^{-6}$ of the payment at
$\nodeshare = 0.20$. Requirement H5 of the mechanism --- targeted recapture below $10^{-4}$ of the
operator share for a single node --- is met with three orders of magnitude to spare. Moreover, the
advantage of \emph{perfect timing} over \emph{random timing} is at most
$1 - q^{\rot{t}} \approx \rot{t}/\drip$ of the long-run dividend, which for one Stratus node at
$\drip = 86{,}400$ is $\SI{1.9}{\percent} \times \SI{0.041}{\percent} \approx 8\times10^{-6}$ of the
operator share.
\end{theorem}

\begin{proof}
$\adv$'s payment is included in block $\hgt-1$, so by \eqref{eq:pnr:inflow}--\eqref{eq:pnr:update}
its contribution $A_\Phi$ is part of $\pool_{\hgt-1}$. By \eqref{eq:pnr:drip} block $\hgt$ drips
$D_\hgt = \lfloor\pool_{\hgt-1}/\drip\rfloor$, of which the tier-$t$ output receives
$\tierratio{t}D_\hgt$ up to the integer residue. The drip is linear in $\pool_{\hgt-1}$, so the part
of $D_\hgt$ attributable to $A_\Phi$ is $A_\Phi/\drip$ and the part reaching $\adv$'s tier-$t$ node
is $\tierratio{t}A_\Phi/\drip$. Substituting $\tierratio{\tS} = 18/27 = 2/3$ and $\drip = 86{,}400$
gives the numeric claim.

For the second statement: $\adv$'s node is paid again at heights $\hgt + j\rot{t}$, $j \ge 1$, and by
the impulse response of Proposition~\ref{prop:pnr:closed} recovers
$\tierratio{t}A_\Phi q^{\,j\rot{t}}/\drip$ each time. Summing,
\[
  \sum_{j \ge 0}\frac{\tierratio{t}A_\Phi q^{\,j\rot{t}}}{\drip}
  = \frac{\tierratio{t}A_\Phi}{\drip\,(1 - q^{\rot{t}})}
  \;\approx\; \frac{\tierratio{t}A_\Phi}{\rot{t}}
  = \frac{\tierratio{t}}{\ncount{t}}A_\Phi ,
\]
using $1 - q^{\rot{t}} \approx \rot{t}/\drip$ for $\rot{t} \ll \drip$ and
$\rot{t} = \ncount{t}$ from Proposition~\ref{prop:pnr:rotation}. This is exactly the un-timed
dividend of Proposition~\ref{prop:pnr:dividend} below. An operator who pays at a uniformly random
phase $U \sim \mathrm{Unif}[0, \rot{t})$ receives the same series delayed, i.e.\ multiplied by
$q^{U} \ge q^{\rot{t}}$. The ratio of best to worst phase is therefore at most
$q^{-\rot{t}}$, and the excess of perfect over random timing is bounded by
$1 - q^{\rot{t}} \approx \rot{t}/\drip = 3{,}246/86{,}400 = \SI{3.76}{\percent}$ of the dividend in
the worst case and \SI{1.9}{\percent} at the Stratus rotation $\rot{\tS} = 1{,}627$.
\end{proof}

The drip does not merely shrink the attack; it removes the value of knowing the schedule.
Knowing the schedule perfectly is worth less than one part in $10^{4}$ of the operator share of a
single payment.

\begin{proposition}[Long-run dividend, \planned]
\label{prop:pnr:dividend}
Because the pool pays out everything it takes in (Property~\ref{prop:pnr:invariants}, I1) in
proportion to the drip, a contributor recovers a fixed share of its own payment over the long run
\emph{regardless of timing}. Define the adversary's drip weight
\[
  w_{\adv} \;\defeq\; \sum_{t \in \tiers} \frac{m_t\,\tierratio{t}}{\ncount{t}}
  \;\in\; [0,1].
\]
Then the long-run recapture of a payment $\price$ is $w_{\adv}\,\nodeshare\,\price$.
\end{proposition}

\begin{proof}
By Proposition~\ref{prop:pnr:closed} every satoshi of $A_\Phi = \nodeshare\price$ is eventually
dripped. By Proposition~\ref{prop:pnr:rotation} node $i$ in tier $t$ receives a fraction
$\tierratio{t}/\ncount{t}$ of all drip in the long run, since it takes $\tierratio{t}$ of the drip
once every $\ncount{t}$ blocks. Summing over $\adv$'s nodes gives $w_{\adv}$, and multiplying by
$A_\Phi$ gives the claim. Since the weight is independent of the height at which the contribution is
made, timing cannot change it.
\end{proof}

\begin{table}[H]
\centering
\small
\begin{tabular}{lrrr}
\toprule
$\adv$ & $w_{\adv}$, of the operator share & of payment, $\nodeshare = 0.20$ & of payment, $\nodeshare = 0.50$\\
\midrule
one Cumulus node   & \SI{0.0023}{\percent} & \SI{0.0005}{\percent} & \SI{0.0011}{\percent}\\
one Nimbus node    & \SI{0.0161}{\percent} & \SI{0.0032}{\percent} & \SI{0.0080}{\percent}\\
one Stratus node   & \SI{0.0410}{\percent} & \SI{0.0082}{\percent} & \SI{0.0205}{\percent}\\
one node of each tier & \SI{0.0593}{\percent} & \SI{0.0119}{\percent} & \SI{0.0297}{\percent}\\
largest measured identity ($237$ Stratus) & \textbf{\SI{9.71}{\percent}} & \textbf{\SI{1.94}{\percent}} & \textbf{\SI{4.86}{\percent}}\\
worst case for the $479$-node identity, if all Stratus & \SI{19.63}{\percent} & \SI{3.93}{\percent} & \SI{9.81}{\percent}\\
\bottomrule
\end{tabular}
\caption{Long-run dividend $w_{\adv}$ by holding. The largest single identity in the measured
census holds $237$ nodes and \SI{10.71}{\percent} of collateral-weighted voting power; the most
nodes held by any one identity is $479$
\srcmeasured{api.runonflux.io/\allowbreak{}daemon/\allowbreak{}listzelnodes}{2026-09-03}.}
\label{tab:pnr:dividend}
\end{table}

The honest reading of Table~\ref{tab:pnr:dividend}: this is not an attack. It is the same fraction
that identity would receive of \emph{everyone else's} payments, it cannot be increased by timing
(Theorem~\ref{thm:pnr:timing}), and it does not distort placement, because placement is not
price-based (\S9). But it is a statement the paper should make plainly rather than leave for a
reader to derive: \textbf{\PNR would give the largest operators a \SIrange{1.9}{4.9}{\percent}
discount on their own hosting}, exactly as \PoN issuance already gives that same identity
\SI{9.71}{\percent} of block rewards \shipped
\srcmeasured{api.runonflux.io/\allowbreak{}daemon/\allowbreak{}listzelnodes}{2026-09-03}. The figure is small because
$\nodeshare \le \nodesharecap = 1/2$ and because the largest identity holds \SI{14.57}{\percent} of
one tier; it would not stay small under further concentration, and the concentration measurements of
\S3 --- a Gini coefficient of $0.8375$ over payment addresses, four identities reaching
\SI{25}{\percent} --- belong next to it.

\begin{proposition}[Wash deployment is a strict loss at every weight, \planned]
\label{prop:pnr:wash}
An adversary paying $\price$ to deploy to itself has net payoff
\[
  \Pi_{\mathrm{wash}} \;=\; -\price\,\big(1 - w_{\adv}\,\nodeshare\big)
  \;\le\; -\tfrac12\,\price
  \qquad\text{for every } w_{\adv} \le 1,\ \nodeshare \le \nodesharecap = \tfrac12 .
\]
\end{proposition}

\begin{proof}
By PS1 the whole of $\price$ leaves the adversary at $\mathsf{sink}$, and by PS4 the Foundation share
$A_F = (1-\nodeshare)\price$ is credited to the Foundation and never returns to it. By
Proposition~\ref{prop:pnr:dividend} it recovers $w_{\adv}\nodeshare\price$ of the pool share and
nothing else. Netting gives the expression; maximising the recovered term over $w_{\adv} \le 1$ and
$\nodeshare \le 1/2$ gives $\Pi_{\mathrm{wash}} \le -\price(1 - 1/2)$.
\end{proof}

\begin{remark}
Even an adversary owning \emph{every node on the network} would lose at least half of every wash
payment at the cap. The largest measured identity would lose \SI{95.14}{\percent} at the cap and
\SI{98.06}{\percent} at activation; a single Stratus operator loses \SI{99.98}{\percent}. The
defence is dilution, and the cap on $\nodeshare$ is what makes it a strict loss at every weight ---
note that the roadmap's original defence, ``rewards are funded only by that workload's own burn,''
presupposed a host-pays model and a burn, and neither exists in the settled design (C1). One
asymmetry survives and should be stated: a \emph{Foundation} workload paying $\price$ returns
$(1-\nodeshare)\price$ to the Foundation, so its effective cost would be $\nodeshare\price$ ---
\SI{20}{\percent} of list price at activation, \SI{50}{\percent} at the cap. Operators are
indifferent, since they receive $\nodeshare\price$ either way, so this is a distortion rather than a
leak: the party that sets the price table would be exposed to one fifth of it.
\end{remark}

% ---------------------------------------------------------------------------
\subsection{\PNR creates no incentive to host}
\label{sec:pnr:hosting}

This is the sharpest question about the design and it is answered here rather than buried.

\begin{property}[Zero marginal hosting incentive, \planned]
\label{prop:pnr:nohosting}
Under Definition~\ref{def:pnr:pool} a node's \PNR income is a function of its tier and its position
in the \PoN queue and of nothing else. In particular
\[
  \frac{\partial\,(\text{node $i$'s \PNR income})}{\partial\,(\text{applications hosted by node } i)}
  \;=\; 0 ,
\]
identically, at every parameter value. The marginal \emph{cost} of hosting --- CPU, memory,
bandwidth, disk, and under \S23 the operational burden of evictions --- is strictly positive.
Therefore the private incentive of every node is to remain in the paid set at minimum real cost.
\end{property}

\begin{proof}
By requirement H4 the \PNR payees of block $\hgt$ are exactly $\payees_\hgt$, and by
Assumption~\ref{as:pnr} A1 $\payees_\hgt$ is determined by the queue keys, which are functions of
last-paid and confirmation heights only. Neither \eqref{eq:pnr:drip} nor \eqref{eq:pnr:tier}
references any placement, metering or workload variable. Hence income is invariant to hosting, and
the derivative is zero.
\end{proof}

\begin{remark}[The collective incentive is real and individually negligible]
\label{rem:pnr:freeride}
Pool inflow depends on aggregate demand, and aggregate demand depends on the network hosting well,
so hosting is a public good among $6{,}489$ nodes. Quantify the private stake. Suppose one node's
refusal to host reduced network revenue by its $1/6{,}489$ share, i.e.\
$\Delta V = 220.7\flux$ per year at $V_0$. That node's own income would fall by
$w\,\nodeshare\,\Delta V$, which at $\nodeshare = 0.20$ is
$0.001\flux$ per year for a Cumulus node, $0.007$ for Nimbus and $0.018$ for Stratus; at
parity-scale revenue ($22\times$) and at the cap, $0.06$--$1.0\flux$ per year. \textbf{Under any
hosting cost above a few satoshis, free-riding dominates for every node at every parameter value in
range.} This is not a tuning problem; it is structural to the settled decision to pay all nodes.
\end{remark}

\begin{remark}[The correct framing, and the cross-section consequence]
\label{rem:pnr:framing}
\PNR is an incentive to \emph{be a node}, not an incentive to \emph{host}. That is a defensible
design --- the network wants capacity available, and availability is what collateral and eligibility
secure --- but it must be said, because a reader who assumes that revenue-funded rewards align
operators with tenants will be wrong.

\textbf{The decoupling is deliberate, and the reason is anti-gaming rather than expediency.} Asked
directly whether some fraction of \PNR should be conditioned on hosting, the team's position is that
it should not: the network decides placement, so paying only the nodes an application landed on
``would be a luck game and can be also misused by node ops that will try to play it trying to force
installation'' \srcintent{08-answers-round-2.md, round 11} (\planned). That reasoning is sound and
follows from a mechanism this paper documents elsewhere. Placement is opportunistic self-selection
with a \SI{90}{\second} broadcast collision wait resolved earliest-broadcast-wins
(\S\ref{sec:fluxos-placement}); attaching income to placement would convert that race into a
financially-motivated one, rewarding whoever is fastest to claim work rather than whoever serves it
best, and giving every operator a standing reason to manipulate what lands on them. Under
opportunistic placement, hosting-conditioned pay would select for latency and manipulation, not for
service quality.

So \Cref{prop:pnr:nohosting} should be read as a \emph{choice with a stated cost}, not as an
oversight. What the choice does not do is remove the cost, and the rest of this remark stands
unchanged: free-riding remains individually rational, and the whole hosting guarantee is transferred
onto the eligibility gate. What actually compels hosting is the \emph{eligibility gate}:
a node is a \PoN payee, and therefore a \PNR payee, only if it is confirmed with a
\fluxbench-signed transaction and, in the target architecture, attested (\S14). \PNR's hosting
guarantee is precisely the strength of that gate and nothing more.

By Assumption~\ref{as:pnr} A4 that gate is weaker than it looks in exactly the relevant way: the
benchmark is self-reported by a process on the operator's own host under network-wide shared keys
(C13), and the attestation signature is not machine-bound (C22). The two findings compose badly.
\textbf{\PNR would raise the payoff to passing the gate without doing the work, by exactly what it
pays} --- \SI{4.29}{\percent} of operator income at activation
(\S\ref{sec:pnr:crossover}), more as the ramp advances (C24). \S4 and \S14 establish the gate's
properties; \S\ref{sec:feasibility} classes the remaining attestation work; and \S\ref{sec:migration} must sequence the attestation fix
\emph{before or with} \PNR activation rather than after it. Attested execution is not adjacent to
\PNR; it is what would make \PNR pay for something. This specification records that as a hard
dependency.
\end{remark}

\begin{openproblem}[Hosting-conditioned eligibility, dependency D1]
\label{op:pnr:d1}
Specify a gate admitting a node to $\payees$ that is conditioned on hosting --- \FluxOS running,
resources real, placed applications actually executing --- and that is bound to a machine rather than
to a shared secret. Interfaces required: a per-node attestation whose signing key is derived from a
hardware root; a verifier reachable by a tenant, not only by loopback mTLS; and a liveness predicate
consumable by the confirmation transaction. The construction is open and belongs to \S14 and \S\ref{sec:feasibility};
without it \PNR pays for membership, and membership is only as real as the gate.
\end{openproblem}

\begin{conjecture}[No hosting incentive is reachable inside the settled constraints]
\label{conj:pnr:nohost}
Within the settled constraints --- beneficiaries are all confirmed nodes (H4), and there is no
metering or attestation dependency between payment and placement (P7) --- no modification of the
settlement rule creates a positive marginal hosting incentive. Payment-conditioned hosting is
excluded by P7; placement is opportunistic self-selection and carries no reputation signal; and no
collateral slashing for non-hosting exists anywhere in the evidence corpus. The conjecture is stated
rather than proved because a proof would require a formal characterisation of the admissible
settlement-rule space, which this specification does not give.
\end{conjecture}

% ---------------------------------------------------------------------------
\subsection{Invariants}
\label{sec:pnr:invariants}

\begin{property}[Settlement invariants, \planned]
\label{prop:pnr:invariants}
Let $\mathrm{In}_\hgt \defeq \sum_{j \le \hgt}F_j$ and
$\mathrm{Paid}_\hgt \defeq \sum_{j \le \hgt}(D_j - \epsilon_j - U_j)$, the drip actually placed in
coinbase outputs. Under Definition~\ref{def:pnr:pool} and Construction~\ref{con:pnr:connect}:
\begin{enumerate}
  \item[\textbf{I1}] \emph{Conservation.} $\pool_\hgt = \mathrm{In}_\hgt - \mathrm{Paid}_\hgt$
    exactly, for all $\hgt \ge \hpa$. Settlement would neither create nor destroy \flux; no term of
    the issuance function $\supply(\hgt)$ of \S5 changes, circulating supply differs from it by
    exactly $\pool_\hgt$ (Property~\ref{prop:pnr:supply}), and there is no burn.
  \item[\textbf{I2}] \emph{Non-negativity.} $\pool_\hgt \ge 0$ for all $\hgt$.
  \item[\textbf{I3}] \emph{The drip never empties the pool.} A positive pool remains positive
    forever, decaying under zero inflow toward a floor below $\drip\sat$
    ($= 0.000864\flux$ at $\drip = 86{,}400$). There is no epoch boundary and no ``last block of the
    pool''.
  \item[\textbf{I4}] \emph{Bounded coinbase.} The coinbase gains no outputs, and the pool commitment
    PC occupies a data field rather than one. Its \emph{operator-facing} value is bounded by
    $\subsidy(\hgt) + \pool_{\hgt-1}/\drip$, and with $\pool_{\hpa-1} = 0$, $F_j \le F_{\max}$ and
    $U_j = 0$ for all $j$, by $\subsidy(\hgt) + F_{\max}$. Its total value additionally carries $G_\hgt$,
    which passes through in one block and is therefore as spiky as demand.
  \item[\textbf{I5}] \emph{Determinism.} $\pool_\hgt$ is a pure function of the block sequence
    through $\hgt$; a disagreement between two implementations surfaces as rejection of the block at
    which it first arises, never as a silent fork.
  \item[\textbf{I6}] \emph{Reorg safety.} Disconnecting block $\hgt$ applies $-\Delta_\hgt$ and
    restores $\pool_{\hgt-1}$ exactly, and $\pool_{\hgt-1}$ is independently readable from the
    parent's coinbase commitment, so recovery needs no journal, no crash marker and no repair RPC.
    On any single chain no payee receives a drip twice for the same block.
  \item[\textbf{I7}] \emph{Liveness independent of demand.} No validity rule references
    $\pool > 0$. $D_\hgt = 0$ is a valid coinbase, identical to today's.
  \item[\textbf{I8}] \emph{No double claim.} A node receives \PNR value only as the increment of its
    \PoN tier output, only in a block in which it is that tier's payee, and at most once per block.
  \item[\textbf{I9}] \emph{Split invariant.} For every output of value $v$ paying $\mathsf{sink}$ in
    a block at height $\hgt$, $A_\Phi \le \nodeshare(\hgt)\,v < A_\Phi + 1\sat$.
\end{enumerate}
\end{property}

\begin{proof}
\textbf{I1} by induction on $\hgt$. Base: $\pool_{\hpa-1} = 0 = \mathrm{In} - \mathrm{Paid}$.
Step: \eqref{eq:pnr:update} gives
$\pool_\hgt = \pool_{\hgt-1} + F_\hgt - (D_\hgt - \epsilon_\hgt - U_\hgt)$, which is exactly the
increment of $\mathrm{In}_\hgt - \mathrm{Paid}_\hgt$. Every satoshi paid to $\mathsf{sink}$ leaves
circulation at payment (PS3) and re-enters through exactly one coinbase output --- $A_F$ through the
dev-fund output, $A_\Phi$ through a tier output as drip over the blocks that follow --- so a
tenant's $v = A_F + A_\Phi$ is fully accounted (PS4). The sum of coinbase values over the chain is
$\sum\subsidy + \sum_j G_j + \mathrm{Paid}_\hgt$, and $\sum_j G_j + \mathrm{Paid}_\hgt \le
\mathrm{Rev}_\hgt$ is value that already existed, so no term of $\supply$ changes;
Property~\ref{prop:pnr:supply} gives the exact circulating-supply identity.

\textbf{I2}: $D_\hgt = \lfloor\pool_{\hgt-1}/\drip\rfloor \le \pool_{\hgt-1}$ for $\drip \ge 1$, and
$\epsilon_\hgt, U_\hgt, F_\hgt \ge 0$, so
$\pool_\hgt \ge \pool_{\hgt-1} - D_\hgt \ge \pool_{\hgt-1}(1 - 1/\drip) \ge 0$.

\textbf{I3}: if $\pool_{\hgt-1} \ge \drip\sat$ then
$\pool_{\hgt-1} - D_\hgt \ge \pool_{\hgt-1}(1 - 1/\drip) > 0$; if $\pool_{\hgt-1} < \drip\sat$ then
$D_\hgt = \lfloor\pool_{\hgt-1}/\drip\rfloor = 0$ and the pool is unchanged. Hence positivity is
preserved in both regimes and the floor is $\drip\sat$.

\textbf{I4}: output count is unchanged because $G_\hgt$ is added to the dev-fund output and
$D_{\hgt,t}$ to the three tier outputs, all four of which already exist, and because PC writes into
a data field rather than creating an output. For the value bound,
$D_\hgt \le \pool_{\hgt-1}/\drip$ is immediate. For the second form, induct: if
$\pool_{\hgt-1} \le \drip F_{\max}$ then
$\pool_\hgt \le \pool_{\hgt-1}(1 - 1/\drip) + \drip F_{\max}/\drip \le \drip F_{\max}$, so
$D_\hgt \le F_{\max}$. Thus the drip never exceeds the largest single-block inflow ever observed,
and $F_{\max}$ is itself bounded by \texttt{MAX\_BLOCK\_SIZE} $= \num{2000000}$~bytes and by
\texttt{MoneyRange}. The $G_\hgt$ term is not smoothed and is not claimed to be: it is a pass-through
to the dev-fund output in the block that receives it, and it reaches no operator, so it affects the
coinbase's size but nothing in \S\ref{sec:pnr:fairness} or \S\ref{sec:pnr:timing}. \emph{Contrast}:
under next-block payout the coinbase would carry the entire \emph{operator} share of a block's
payments and would spike with demand, which is the property that would break fairness.

\textbf{I5}: $\pool_\hgt$ is defined by integer arithmetic on the block sequence, so any node with
the chain recomputes it. Two implementations that disagree at $\hgt$ disagree with the commitment PC
carried by that block (step 6 of Construction~\ref{con:pnr:connect}), so the block at which the
divergence arises is the block that is rejected --- a detected fault, localised, and one block
earlier than a per-index delta would have surfaced it.

\textbf{I6}: $\Delta_\hgt$ is subtracted on disconnect, which is exact because
\eqref{eq:pnr:update} is additive; and because $\pool_{\hgt-1}$ is committed in the parent's
coinbase, correctness does not depend on the cache surviving. A \PNR application payment in an
orphaned block returns to the mempool and re-enters the pool when re-mined; the orphaned coinbase
and its drip vanish with the block. \emph{Proof sketch only for the implementation obligation}: the
pool must not be a global mutable singleton. The documented reorg fragility of
\texttt{g\_fluxnode\allowbreak{}Cache} is the anti-pattern to avoid
\srccode{fluxd/\allowbreak{}src/\allowbreak{}fluxnode/\allowbreak{}fluxnode.cpp}{790--820}, and the in-tree precedent to follow is the
per-index shielded value-pool accumulator (\S20).

\textbf{I7}: no clause of Construction~\ref{con:pnr:connect} tests $\pool > 0$; step 2 yields
$D_\hgt = 0$ when $\pool_{\hgt-1} < \drip$, and step 4 then requires the tier outputs to equal
$\subsidy_t(\hgt)$, which is today's rule.

\textbf{I8}: the payee set is $\payees_\hgt$, one node per tier (Assumption~\ref{as:pnr} A1), and the
drip is added to that tier's single output. A node whose collateral moves re-enters the queue at the
back with last-paid reset and waits one full rotation; it carries no accumulated claim, so there is
nothing to double-claim.

\textbf{I9}: PS4 sets $A_\Phi = \lfloor r(\hgt)v/100\rfloor$ with
$r(\hgt) = 100\nodeshare(\hgt)$ an integer, and $\lfloor x \rfloor \le x < \lfloor x \rfloor + 1$.
\end{proof}

\begin{property}[Two security constraints the design exports, \planned]
\label{prop:pnr:exported}
The following two rules are not optimisations. Each is a defect if violated.
\begin{enumerate}
  \item[\textbf{S1}] \emph{Emergency blocks must not drip.} For any block produced through the
    2-of-4 emergency key path (Assumption~\ref{as:pnr} A3), $D_\hgt \defeq 0$ and $\pool_\hgt
    \defeq \pool_{\hgt-1} + F_\hgt$.
  \item[\textbf{S2}] \emph{Moderation forfeitures must never enter $\pool$.} If \S23 forfeits a
    removed tenant's prepaid balance, that balance must go anywhere other than the fee pool.
\end{enumerate}
\end{property}

\begin{proof}[Justification]
\textbf{S1}. An emergency block is produced outside the \PoN eligibility check, with no time or
stall gating \srccode{fluxd/\allowbreak{}src/\allowbreak{}emergencyblock.cpp}{183--191}, and the
remainder output of its coinbase defaults to the dev-fund address, as in every block
\srccode{fluxd/\allowbreak{}src/\allowbreak{}rpc/\allowbreak{}emergencyblock.cpp}{64--77},
\srccode{fluxd/\allowbreak{}src/\allowbreak{}pon/\allowbreak{}pon-minter.cpp}{288--306} (C9). Its
coinbase is nonetheless bound by the deterministic payout check, so its tier outputs go to the
queue's payees exactly as in any other block
\srccode{fluxd/\allowbreak{}src/\allowbreak{}miner.cpp}{487--489},
\srccode{fluxd/\allowbreak{}src/\allowbreak{}main.cpp}{3878--3884}. The capability the drip would
add to the emergency path is therefore not redirection but acceleration: the holders of two of four
hardcoded keys could produce blocks faster than the \SI{30}{\second} target and run the drip
forward, $\pool_{\hgt-1}/\drip$ per block --- at the recommended $\drip$ and parity-scale revenue,
$15.2\flux$ per block from a pool holding $1.31$M$\flux$ --- paid to whichever nodes stand at the
front of the queue, including any they operate, with no gating on how many such blocks are
produced. Setting $D_\hgt = 0$ removes that lever entirely while preserving I1 (nothing is paid, so
$\mathrm{Paid}$ does not advance) and I7 (a zero drip is a valid coinbase). The Foundation term
$G_\hgt$ needs no equivalent rule and is left to flow: an emergency block that contains sink
payments would pay their \SI{80}{\percent} to the dev-fund address, which is the destination that
rule already has, so the emergency path gains no capability it did not have. The drip is the only
term over which it would gain one, and S1 removes exactly that.

\textbf{S2}. Routing forfeitures into $\pool$ would give every node a positive expected payoff from
every eviction, since by Proposition~\ref{prop:pnr:dividend} each node receives
$\tierratio{t}/\ncount{t}$ of everything the pool takes in. Report brigading would then be paid for
by the settlement mechanism, and the moderation quorum of \S23 --- which is Sybil-resistant by
construction but, at the measured distribution, not censorship-resistant (C6) --- would acquire a
direct financial motive. The constraint is a requirement exported to \S23, not a preference; note
also that ``no burn'' is a \PNR decision and does not bind \S23's choice of destination.
\end{proof}

\settlednote{$D_\hgt = 0$ on emergency blocks (\Cref{tab:pnr-recommended}). O15 --- emergency blocks and the drip. Domain: $\{D_\hgt = 0\ \text{(required by S1)},\
\text{normal drip}\}$. Originally recorded as open because it is a consensus rule nobody has written; S1 states
what it must be.}

\begin{remark}[O16, destination of forfeited balances --- burn, and why the question matters]
\label{rem:o16-burn}
This is the one open parameter in this section whose \emph{wrong} answer creates an attack rather
than an inefficiency, so it is worth restating why it was asked before recording the recommendation.

Two forfeiture events exist and they are not the same. \textbf{Ordinary expiry} forfeits nothing: an
application's prepaid period is consumed as it runs, the application expires when it is exhausted,
its data is deleted, and there is no refund and no residue \srcintent{08-answers-round-2.md, round
5}. \textbf{Moderation eviction} is different, because it interrupts a period that has been paid for
and not yet consumed: an application evicted at height $\hgt_b$ with paid-through height
$e_\app > \hgt_b$ leaves an unearned remainder, and the author has settled that this remainder is
forfeited rather than refunded \srcintent{08-answers-round-2.md, round 5}. The question is where that
value goes.

\textbf{It must not go to $\pool$.} If forfeited balances flowed to the fee pool, the operators who
vote to evict an application would be paid, pro rata, out of the balance they voted to forfeit.
Under the \SI{25}{\percent} collateral-weighted threshold of \S\ref{sec:governance}, and given the
concentration measured there, that is a direct financial incentive to evict paying applications ---
strongest precisely against the largest prepaid balances, which are the most committed tenants. This
is why S2 forbids it, and it is the substance of the question rather than an accounting detail.

\textbf{Recommendation: burn.} Of the three admissible destinations, burning is the only one that
leaves no party with a financial interest in the outcome of a vote. A refund to the tenant is
excluded by the settled forfeiture decision. Routing to the Foundation is defensible --- the
Foundation does not vote, so the incentive is not created at the voter --- but it makes the
Foundation a beneficiary of network moderation, which is an association worth paying a small price to
avoid, and the amounts are small enough that the price is small: the enforcement floor is
\num{0.01}\flux\ per application (\S\ref{sec:governance}). Burning also composes correctly with the
supply argument of \S\ref{sec:monetary-policy}, since a burn is exactly representable in the supply
function whereas a Foundation transfer is not.

\settlednote{O16: \textbf{burn} the unearned remainder on moderation eviction, adopted
\srcintent{08-answers-round-2.md, round 11}. $\pool$ is excluded by S2, and refund by the settled
forfeiture position. The implementing choice --- an explicit output to a provably unspendable script,
or an omission from the coinbase --- belongs to \S\ref{sec:governance}, which owns the enforcement
path; both realise the same supply effect.}
\end{remark}

\paragraph{Failure modes not covered by the invariants.}
Three deserve naming. \emph{A scheduled payee that is offline at payout} resolves under the existing
rule: either it is still confirmed, in which case it is paid --- drip included, exactly as its
subsidy is --- or it has failed to re-confirm within $640$ blocks and is purged before payout
\shipped \srccode{fluxd/\allowbreak{}src/\allowbreak{}fluxnode/\allowbreak{}fluxnode.cpp}{790--820}. The consequence a reader should see is
that the drip, like the subsidy, would be paid for being in the set, not for being up at the block.
\emph{Sustained low demand} decays the pool geometrically to dust (I3) and the coinbase to today's,
with no liveness effect (I7); since \PNR would be \SI{4.29}{\percent} of operator income at
activation, low demand is the baseline rather than a failure mode. \emph{A deep reorg} --- normally
bounded at $40$ blocks, but temporarily $5{,}000$ during the \PoN stabilization window
\srccode{fluxd/\allowbreak{}src/\allowbreak{}main.cpp}{8983--8988} (C16) --- undoes drips and inflows together and is tolerated
by I6. The single-output shape removes an obligation that the two-output shape carried: because the
payer commits to no split, a payment returned to the mempool cannot become invalid at a new height.
It would simply be divided at the height at which it is re-mined, so a payment straddling a
reduction boundary changes the destination of $\rampstep = 0.05$ of its value and nothing else,
and \FluxOS's coverage check is unaffected.

% ---------------------------------------------------------------------------
\subsection{Activation schedule}
\label{sec:pnr:activation}

\begin{table}[H]
\centering
\small
\begin{tabular}{lrlrl}
\toprule
event & height & date & $\subsidy$ (\flux/block) & $\nodeshare$\\
\midrule
current tip (baseline, \shipped)      & $\sim$2{,}912{,}015 & 2026-09-03 & $14.0$   & ---\\
reduction 1 ($\times 9/10$)           & 3{,}071{,}200 & 2026-10-25 & $12.6$   & ---\\
\textbf{PA Depletion} ($\times 1/2$)  & \textbf{3{,}466{,}630} & \textbf{2027-03-12} & $\mathbf{6.3}$ & $\mathbf{0.20}$ --- \PNR begins\\
reduction 2                            & 4{,}122{,}400 & 2027-10-25 & $5.67$   & $0.25$\\
reduction 3                            & 5{,}173{,}600 & 2028-10-24 & $5.103$  & $0.30$\\
reduction 4                            & 6{,}224{,}800 & 2029-10-24 & $4.5927$ & $0.35$\\
reduction 5                            & 7{,}276{,}000 & 2030-10-24 & $4.1334$ & $0.40$\\
reduction 6                            & 8{,}327{,}200 & 2031-10-24 & $3.7201$ & $0.45$\\
reduction 7 (cap reached)              & 9{,}378{,}400 & 2032-10-23 & $3.3481$ & $\mathbf{0.50} = \nodesharecap$\\
annually thereafter                    & $+\redint$    &            & $\times 9/10$, $20$ times & $0.50$\\
\bottomrule
\end{tabular}
\caption{\PNR activation and ramp \planned. Heights lie exactly on the \PoN reduction grid
$\hpon + k\,\redint$ with $\hpon = 2{,}020{,}000$ and $\redint = 1{,}051{,}200$
\srccode{fluxd/\allowbreak{}src/\allowbreak{}chainparams.cpp}{153,157}; the subsidy column is the as-coded per-step-truncating
path $14\cdot(9/10)^{\mathrm{red}(\hgt)}$ multiplied by $1/2$ from $\hpa$
\srcintent{evidence/design/02-pnr.md}. Dates are $\SI{30}{\second}$ extrapolations from the \PoN
activation timestamp anchor and carry the drift documented in \S\ref{sec:evaluation} (about $\SI{0.8}{\day}$ over
$796{,}820$ blocks); the intake gives one of them a day later. The intake says the cap is reached
``$\sim$2033''; the grid gives October 2032.}
\label{tab:pnr:activation}
\end{table}

\settlednote{Uniform $\times 1/2$ (\Cref{tab:pnr-recommended}). O20 --- whether the PA-Depletion halving applies uniformly to the three tier bases and to
the dev-fund remainder. Domain: $\{\text{uniform } \times 1/2,\ \text{other}\}$.
Table~\ref{tab:pnr:activation} and every parity figure in \S\ref{sec:pnr:crossover} assume uniform;
a non-uniform cut changes operator issuance and therefore every crossover number.}

\settlednote{Fixed schedule plus a stated review (\Cref{tab:pnr-recommended}); see also \Cref{rem:pnr:contraction} on the contraction response. O13 --- ramp policy if demand undershoots. Domain: $\{\text{fixed schedule (settled as far
as it goes)},\ \text{a revision process}\}$. With $g = 0$ there is no crossover ever
(\S\ref{sec:pnr:crossover}); the schedule's behaviour in that case is unstated.}

\begin{remark}[Sequencing conflict C1a --- settled as simultaneous]
\label{rem:c1a-settled}
The design intake activated \PNR{} at PA~Depletion simultaneously with the parallel-asset emission
cut, while the public roadmap stated that parallel-asset mining retirement comes \emph{only after}
\PNR{} is demonstrably paying operators, ninety days announced, with the two in different halves of
2027. The team has settled it: activation is \textbf{simultaneous}, both at
$\hpa = \num{3466630}$ \srcintent{08-answers-round-2.md, round 5} (\planned), superseding the
roadmap's strictly-after ordering (\S\ref{sec:pa-sequencing}).

The consequence is recorded rather than smoothed over, because simultaneity removes a check the
roadmap's ordering provided: \textbf{no external observer can confirm that \PNR{} pays operators
before parallel-asset mining rewards are retired}, since both happen at one height. The roadmap's
sequence would have produced that evidence; the settled sequence does not. \S\ref{sec:end-to-end} carries
this as an accepted risk that should stay visible in any activation announcement rather than as an
oversight.
\end{remark}

% ---------------------------------------------------------------------------
\subsection{The crossover, in both denominations}
\label{sec:pnr:crossover}

\paragraph{Denomination discipline.} \flux terms are primary, because that is what consensus
governs. Real terms are a clearly-labelled secondary analysis with price treated as an exogenous
parameter. No price prediction, target, or investment framing appears anywhere in this section, and
a claim made in one denomination is never defended with reasoning from the other
\srcintent{evidence/design/07-pricing-and-price.md}. \textbf{Parity depends on demand growth, not on
price.}

\paragraph{Bases.} Operator issuance excludes the dev-fund remainder: it is
$\tfrac{13.5}{14}\,\subsidy\,\blocksyear$, not $\subsidy\,\blocksyear$, because $0.5/14$ of every
\PoN subsidy is a consensus-enforced dev-fund payment \shipped
\srccode{fluxd/\allowbreak{}src/\allowbreak{}main.cpp}{2877--2884} (C9). \textbf{The intake's own parity table is
\SI{3.7}{\percent} high} for exactly this reason: it reports $33{,}112{,}800\flux$ per year of parity
revenue at activation where the operator-only basis gives $31{,}930{,}200$, and
$6{,}622{,}560\flux$ per year of issuance where the operator-only basis gives $6{,}386{,}040$
\srcintent{evidence/design/02-pnr.md}. The generated data carries both bases
\srcdoc{paper/data/pnr\_crossover.dat}.

\paragraph{At activation.} With $\nodeshare = 0.20$ and $\subsidy = 6.3\flux$ per block, operator
issuance would be $6{,}386{,}040\flux$ per year and \PNR income at $V_0$ would be
$0.20 \times 1{,}431{,}600 = 286{,}320\flux$ per year. Defining the \PNR share of operator income as
$\text{\PNR}/(\text{issuance} + \text{\PNR})$,
\[
  \frac{286{,}320}{6{,}386{,}040 + 286{,}320} \;=\; \SI{4.29}{\percent}.
\]
Parity revenue --- the annual application revenue at which \PNR income equals operator issuance ---
would be $6{,}386{,}040/0.20 = 31{,}930{,}200\flux$ per year, or $22.3\times V_0$.

\subsubsection{\flux terms (primary)}
\label{sec:pnr:fluxterms}

\begin{table}[H]
\centering
\small
\begin{tabular}{rlrrrrr}
\toprule
stage start & date & $\subsidy$ & $\nodeshare$ & operator issuance & parity revenue & growth needed\\
(height) & & (\flux) & & (\flux/yr) & (\flux/yr) & from $V_0$\\
\midrule
3{,}466{,}630  & 2027-03-12 & $6.3000$ & $0.20$ & 6{,}386{,}040 & 31{,}930{,}200 & ---\\
4{,}122{,}400  & 2027-10-25 & $5.6700$ & $0.25$ & 5{,}747{,}436 & 22{,}989{,}744 & $16.1\times$\\
5{,}173{,}600  & 2028-10-24 & $5.1030$ & $0.30$ & 5{,}172{,}692 & 17{,}242{,}308 & $12.0\times$\\
6{,}224{,}800  & 2029-10-24 & $4.5927$ & $0.35$ & 4{,}655{,}423 & 13{,}301{,}209 & $9.3\times$\\
7{,}276{,}000  & 2030-10-24 & $4.1334$ & $0.40$ & 4{,}189{,}881 & 10{,}474{,}702 & $7.3\times$\\
8{,}327{,}200  & 2031-10-24 & $3.7201$ & $0.45$ & 3{,}770{,}893 & 8{,}379{,}762  & $5.9\times$\\
9{,}378{,}400  & 2032-10-23 & $3.3481$ & $\mathbf{0.50}$ & 3{,}393{,}803 & 6{,}787{,}607 & $4.7\times$\\
12{,}532{,}000 & 2035-10-23 & $2.4407$ & $0.50$ & 2{,}474{,}083 & 4{,}948{,}165  & $3.5\times$\\
17{,}788{,}000 & 2040-10-21 & $1.4412$ & $0.50$ & 1{,}460{,}921 & 2{,}921{,}842  & $2.0\times$\\
23{,}044{,}000 & 2045-10-20 & $0.8510$ & $0.50$ & 862{,}659     & 1{,}725{,}319  & $1.2\times$ (perpetual tail)\\
\bottomrule
\end{tabular}
\caption{Parity requirement in \flux terms \planned. The two levers compound as intended: the
requirement falls from $31.9$M to $6.8$M\flux\ per year in $5.6$ years and to $1.7$M on the
perpetual tail. The final row is the flat tail of \S5: after twenty reductions the subsidy freezes,
so operator issuance never reaches zero \srccode{fluxd/\allowbreak{}src/\allowbreak{}chainparams.cpp}{158}.
Source: \texttt{paper/\allowbreak{}data/\allowbreak{}pnr\_crossover.\allowbreak{}dat}.}
\label{tab:pnr:parity}
\end{table}

\begin{table}[H]
\centering
\small
\begin{tabular}{lllr}
\toprule
annual revenue growth $g$ & crossover & revenue at crossover (\flux/yr) & $\nodeshare$ at crossover\\
\midrule
\SI{0}{\percent}   & \textbf{never} & --- & ---\\
\SI{10}{\percent}  & 2038-10 & $4.33$M & $0.50$\\
\SI{25}{\percent}  & 2033-10 & $6.28$M & $0.50$\\
\SI{50}{\percent}  & 2031-10 & $9.33$M & $0.45$\\
\SI{100}{\percent} & 2030-10 & $17.6$M & $0.40$\\
\SI{200}{\percent} & 2029-10 & $25.6$M & $0.35$\\
\bottomrule
\end{tabular}
\caption{Crossover year under constant annual revenue growth from $V_0$, \flux terms \planned.
Required sustained growth to reach parity by a given year: 2029, \SI{134}{\percent};
2030, \SI{73}{\percent}; 2031, \SI{47}{\percent}; 2032, \SI{32}{\percent}; 2033, \SI{25}{\percent};
2035, \SI{16}{\percent}. Source: \texttt{paper/\allowbreak{}data/\allowbreak{}pnr\_crossover.\allowbreak{}dat}, columns
\texttt{pnr\_g000}\,\dots\,\texttt{pnr\_g200}.}
\label{tab:pnr:crossover}
\end{table}

\textbf{With flat demand there is no crossover, ever.} The perpetual tail issuance of
$862{,}659\flux$ per year to operators exceeds $\nodesharecap \cdot V_0 = 0.5 \times 1{,}431{,}600 =
715{,}800\flux$ per year, and the tail is perpetual because the subsidy freezes after twenty
reductions rather than going to zero (\S5, C7). So the correct statement, in these words: \emph{as
scheduled, \PNR would begin as a supplement; whether it becomes the main component of operator income
is a statement about demand, and this paper makes no such statement.}

\begin{remark}[The stated response to a demand shortfall is contraction, and it works]
\label{rem:pnr:contraction}
Asked what happens if demand does not arrive, the team's answer is that the network shrinks: hosting
that many FluxNodes stops being economic, some operators leave, ``which would mean more flux
allocation for other node ops'' \srcintent{08-answers-round-2.md, round 11} (\planned). The team
also reports the observed direction is the opposite --- demand rising, with faster growth
anticipated from AI-driven deployments --- and this paper records that as stated intent, not as a
measurement it can confirm.

The equilibrium argument is sound and worth making precise, because it is the mechanism's genuine
answer to its own worst case. \PoN{} pays a fixed per-block subsidy divided among confirmed nodes,
so per-node issuance income is inversely proportional to the confirmed node count. If total operator
income falls below the cost of operating a node, exit raises the income of every node that remains,
and exit continues until income again covers cost. The node count is therefore self-correcting
against demand: the network converges on the size its revenue plus issuance can support, with no
governance intervention and no parameter change. That is a stronger position than most
demand-funded designs hold, and it does not depend on demand materialising.

\textbf{The cost is a node count that falls exactly when the paper's other mechanisms need it to
rise}, and the two answers are in tension in a way neither alone reveals.
\Cref{sec:23:measurement} measures four owner identities reaching the \SI{25}{\percent}
moderation threshold, and the stated remedy for that concentration is growth: ``as network grows it
will be more decentralized'' \srcintent{08-answers-round-2.md, round 11}. But the stated response to
weak demand is contraction --- and contraction concentrates. Under the demand shortfall this
subsection describes, the operators who exit first are the marginal ones, which are disproportionately
single-node operators rather than the large holders whose collateral drives the concentration
measurement. The concentration ratio would therefore \emph{worsen} in precisely the scenario where
the mechanism relies on it improving.

Both answers are individually reasonable and the paper endorses neither as a prediction. What it
records is the dependency: \textbf{the moderation quorum's censorship-resistance is a function of
demand for Flux Cloud}, through node count, and no part of the moderation design references demand at
all. \Cref{sec:23:remedies} evaluates remedies on the measured distribution; none of them is robust
to that distribution getting worse.
\notestablished{the composition of exit under a demand shortfall --- whether marginal
operators are disproportionately single-node holders --- is not measured in this corpus. The
directional claim that contraction concentrates follows from the tier and collateral structure, but
its magnitude is not established.}
\end{remark}

\begin{figure}[H]
\centering
\begin{tikzpicture}
\begin{axis}[
  width=0.92\linewidth, height=6.4cm,
  xlabel={year},
  ylabel={\flux\ per year (log scale)},
  ymode=log,
  xmin=2027, xmax=2046,
  legend pos=north west,
  legend cell align=left,
  legend style={font=\scriptsize},
  grid=major,
  every axis plot/.append style={thick},
]
\addplot[black,mark=*,mark size=1.2pt] coordinates {
  (2027.19,6386040) (2027.81,5747436) (2028.81,5172692) (2029.81,4655423)
  (2030.81,4189881) (2031.81,3770893) (2032.81,3393803) (2035.81,2474083)
  (2040.80,1460921) (2045.80,862659)};
\addlegendentry{operator issuance, $\tfrac{13.5}{14}\subsidy\blocksyear$}
\addplot[blue,densely dashed,mark=square*,mark size=1.1pt] coordinates {
  (2027.19,286320) (2027.81,357900) (2028.81,429480) (2029.81,501060)
  (2030.81,572640) (2031.81,644220) (2032.81,715800) (2035.81,715800)
  (2040.80,715800) (2045.80,715800)};
\addlegendentry{\PNR income, $g=\SI{0}{\percent}$}
\addplot[teal,dashdotted,mark=triangle*,mark size=1.3pt] coordinates {
  (2027.19,286320) (2027.81,411010) (2028.81,616517) (2029.81,899102)
  (2030.81,1284395) (2031.81,1806178) (2032.81,2508521) (2035.81,4899436)
  (2040.80,14918700) (2045.80,45546000)};
\addlegendentry{\PNR income, $g=\SI{25}{\percent}$}
\addplot[red,dotted,mark=diamond*,mark size=1.3pt] coordinates {
  (2027.19,286320) (2027.81,550090) (2028.81,1320222) (2029.81,3081018)
  (2030.81,7041655) (2031.81,15844590) (2032.81,35210000)};
\addlegendentry{\PNR income, $g=\SI{100}{\percent}$}
\end{axis}
\end{tikzpicture}
\caption{Crossover in \flux terms \planned. Operator issuance falls on the reduction grid while the
operator share $\nodeshare$ ramps from $0.20$ to the cap $\nodesharecap = 0.50$; \PNR income is
$\nodeshare(\hgt)\,V_0(1+g)^{t}$. The $g = \SI{0}{\percent}$ series flattens at
$715{,}800\flux$ per year and never meets the issuance curve, whose perpetual tail is
$862{,}659\flux$ per year. Price does not appear on either axis and no price claim is made.}
\label{fig:pnr:crossover}
\end{figure}

\begin{remark}[Figure provenance]
The coordinates plotted in Figures~\ref{fig:pnr:crossover} and~\ref{fig:pnr:real} are the tabulated
output of \texttt{scripts/\allowbreak{}07-pnr-settlement.\allowbreak{}py} (run 2026-09-03), which also wrote
\texttt{paper/\allowbreak{}data/\allowbreak{}pnr\_crossover.\allowbreak{}dat} and \texttt{paper/\allowbreak{}data/\allowbreak{}pnr\_crossover\_real.\allowbreak{}dat}. Every
plotted point in Figure~\ref{fig:pnr:crossover} has been checked against the corresponding row of
that data file and agrees exactly; the data files carry their own \texttt{\#} provenance headers and
the figures are reproducible from them.
\end{remark}

\subsubsection{Real terms (secondary, price exogenous)}
\label{sec:pnr:real}

The \flux-denominated price table is revised by policy, and the revisions have historically tracked
real cost: the CPU rate has fallen $100\times$ across five revisions \shipped
\srcmeasured{api.runonflux.io/\allowbreak{}apps/\allowbreak{}deploymentinformation}{2026-09-03}. Under a continuation of that
rebasing policy, \flux inflow to the pool would be real demand divided by an exogenous price multiple,
$V^{\flux}_t = V_0(1+g)^{t}/m_t$, while issuance is fixed in \flux. Treating $m$ as an axis and not
as a forecast gives the following years-to-parity after activation.

\begin{table}[H]
\centering
\small
\begin{tabular}{lrrrrr}
\toprule
real growth $g$ $\backslash$ price multiple $m$ & $0.5$ & $1$ & $2$ & $5$ & $10$\\
\midrule
\SI{0}{\percent}   & $14.6$ & never & never & never & never\\
\SI{10}{\percent}  & $7.6$  & $11.6$ & $14.6$ & $19.6$ & $26.1$\\
\SI{25}{\percent}  & $5.6$  & $6.6$  & $9.6$  & $11.6$ & $13.6$\\
\SI{50}{\percent}  & $3.6$  & $4.6$  & $5.6$  & $7.6$  & $9.6$\\
\SI{100}{\percent} & $2.6$  & $3.6$  & $4.6$  & $5.6$  & $5.6$\\
\bottomrule
\end{tabular}
\caption{Years from activation to \flux-terms parity, as a function of real demand growth $g$ and an
exogenous price multiple $m$ \planned. $m$ is a sensitivity axis, not a projection.
Source: \texttt{paper/\allowbreak{}data/\allowbreak{}pnr\_crossover\_real.\allowbreak{}dat}.}
\label{tab:pnr:real}
\end{table}

\begin{figure}[H]
\centering
\begin{tikzpicture}
\begin{axis}[
  width=0.74\linewidth, height=6.0cm,
  xlabel={exogenous price multiple $m$ (log scale)},
  ylabel={years from activation to \flux-terms parity},
  xmode=log,
  log basis x=10,
  xmin=0.4, xmax=12,
  ymin=0, ymax=28,
  legend cell align=left,
  %% Outside the plot box: at north west the g=10% series ran straight through it.
  legend style={at={(1.03,1)}, anchor=north west, font=\scriptsize, draw=black!30},
  grid=major,
  every axis plot/.append style={thick},
]
\addplot[blue,mark=*,mark size=1.2pt] coordinates {
  (0.5,7.6) (1,11.6) (2,14.6) (5,19.6) (10,26.1)};
\addlegendentry{$g=\SI{10}{\percent}$}
\addplot[teal,densely dashed,mark=square*,mark size=1.1pt] coordinates {
  (0.5,5.6) (1,6.6) (2,9.6) (5,11.6) (10,13.6)};
\addlegendentry{$g=\SI{25}{\percent}$}
\addplot[orange,dashdotted,mark=triangle*,mark size=1.3pt] coordinates {
  (0.5,3.6) (1,4.6) (2,5.6) (5,7.6) (10,9.6)};
\addlegendentry{$g=\SI{50}{\percent}$}
\addplot[red,dotted,mark=diamond*,mark size=1.3pt] coordinates {
  (0.5,2.6) (1,3.6) (2,4.6) (5,5.6) (10,5.6)};
\addlegendentry{$g=\SI{100}{\percent}$}
\addplot[black,only marks,mark=x,mark size=2.5pt] coordinates {(0.5,14.6)};
\addlegendentry{$g=\SI{0}{\percent}$ (parity unreachable for $m \ge 1$)}
\end{axis}
\end{tikzpicture}
\caption{Real-terms sensitivity \planned, secondary analysis. Reading, and it is the point of the
figure: under a rebasing price policy, \emph{price appreciation delays \flux-terms parity} --- every
curve rises to the right. Appreciation can raise operator income in real terms while the
\flux-denominated ratio barely moves. The two statements live in different denominations and must
never be used to defend one another.}
\label{fig:pnr:real}
\end{figure}

\settlednote{Trailing receipts at $\mathsf{sink}$ (\Cref{tab:pnr-recommended}). O17 --- the measured revenue baseline that should replace ``$100\flux$ per application per
month''. Domain: $\{\text{trailing receipts at the payment sink } \texttt{t3Nryf\ldots}$,
$\text{list-price value of the measured resource totals}\}$; the census gives $8{,}972$ vCPU,
$17{,}544{,}229$~MB of RAM and $160{,}107$~GB of disk across $7{,}486$ instances
\srcmeasured{api.runonflux.io/\allowbreak{}apps/\allowbreak{}globalappsspecifications}{2026-09-03}. Every absolute figure in
this subsection scales linearly with it; the growth-rate conclusions do not.}

\settlednote{Rule-based rebasing (\Cref{tab:pnr-recommended}). The
decision is that rebasing follows a published rule rather than discretion or an oracle; the rule's
functional form is specified in \S\ref{sec:pnr:recommended} as trailing measured receipts adjusted at
fixed intervals, and its coefficients are a policy calibration rather than a mechanism choice. O9 --- price-table rebasing policy. Domain: $\{\text{discretionary (today)},\
\text{rule-based},\ \text{USD oracle}\}$. A rule removes a Foundation lever over pool inflow; an
oracle adds a trust dependency that would need its own treatment. The existence of a
\texttt{getappspecsusdprice} work item suggests movement here.}

% ---------------------------------------------------------------------------
\subsection{Open parameters}
\label{sec:pnr:openparams}

Of the twenty decisions this specification opened, four are now settled and are recorded in
Table~\ref{tab:pnr:closed}; the settlement decision created one more, O21, which \Cref{rem:pnr:o21} settles. Eleven more are
settled by the author's confirmation of \Cref{tab:pnr-recommended}, and O16 by \Cref{rem:o16-burn}; four therefore
remain open (O4, O8, O14 and O19), and the load-bearing one among them is \textbf{O8} --- which payment channels must
settle through $\mathsf{sink}$ --- because every channel outside it is revenue the consensus split
never touches and the counter of Property~\ref{prop:pnr:supply} never sees. None of the open values
is guessed anywhere in this section, and none of the closed ones is stated without its provenance.

\begin{table}[H]
\centering
\scriptsize
\begin{tabular}{llp{0.42\linewidth}l}
\toprule
id & decision & settled value & provenance\\
\midrule
O1 & drip constant $\drip$ & $86{,}400$ blocks, i.e.\ exactly \SI{30}{\day} at \SI{30}{\second} & \srcintent{evidence/design/08-answers-round-2.md}\\
O2 & commitment of $\pool_\hgt$ & committed in the coinbase (PC) and recomputed by every validator; no separate store, no header field & \srcintent{evidence/design/08-answers-round-2.md}, \srcdoc{plan/conflicts.md C37}\\
O3 & memo semantics & \texttt{OP\_RETURN} binding $H(m_\app)$, carried and not interpreted by consensus (PS2) & \srcintent{evidence/design/08-answers-round-2.md}\\
O18 & serialisation and enforcement boundary & one output of the full price to $\mathsf{sink}$ (PS1); consensus applies the split (PS4); \FluxOS checks only existence, memo and coverage & \srcintent{evidence/design/08-answers-round-2.md}, \srcdoc{plan/conflicts.md C36, C38}\\
\bottomrule
\end{tabular}
\caption{Decisions closed since the first draft of this specification \planned. All four are team
ratifications: O2 originated as this specification's recommendation and was subsequently adopted.
Each would be a consensus rule once activated, so changing any of them afterwards would be a network
upgrade.}
\label{tab:pnr:closed}
\end{table}

\begin{table}[H]
\centering
\scriptsize
\begin{tabular}{llp{0.30\linewidth}p{0.34\linewidth}}
\toprule
id & decision & domain & trade-off\\
\midrule
O4 & fee floor and spam resistance & $p_{\min} \ge 0.01\flux$; per-instance or per-deployment & dust-priced deployments contribute nothing yet occupy capacity; a floor prices out the micro-workloads \S18 wants\\
O5 & $D_{\hgt,t}$ with no payee in tier $t$ & stays in pool; dev fund; other tiers & pool: conservation-trivial; dev fund: one rule for subsidy and drip; other tiers: rewards concentration\\
O6 & partition residue $\epsilon_\hgt$ & pool; dev fund & $\le 2\sat$ per block either way; must be fixed to avoid divergence\\
O7 & shielded payers & transparent value required; Sapling with revealed amount; excluded & consensus needs the received value for PS4; the hide-payer/reveal-amount construction is \S21's\\
\textbf{O8} & scope of payments required to settle at $\mathsf{sink}$ & new deployments; renewals; credit top-ups (\S16); \CumulusVPN; enterprise & every excluded channel is revenue outside the consensus split and invisible to the revenue counter\\
O9 & price-table rebasing & discretionary; rule-based; USD oracle & a rule removes a Foundation lever; an oracle adds a trust dependency\\
O10 & drip ordering & from $\pool_{\hgt-1}$; after applying $F_\hgt$ & parent-state: coinbase checkable early, producer cannot influence own drip; post-inflow: one block less lag, same-block recapture returns\\
O11 & initial seed $\pool_{\hpa-1}$ & $0$; Foundation-funded seed & seed removes the $\sim$90-day ramp but is a one-off transfer and must be labelled one\\
O12 & fleet-growth guard on $\drip$ & none; documented hard-fork review trigger & at $4\times$ Cumulus the spread is \SI{16}{\percent}; an automatic rule is a manipulation surface\\
O13 & ramp policy if demand undershoots & fixed schedule; a stated revision process & with $g=0$ there is no crossover; the schedule's behaviour is unstated\\
O14 & identity of $\mathsf{sink}$ & keep \texttt{t3Nryf\ldots}; a fresh address; rotation/PQ path (\S8) & keeping it preserves tooling and one continuous history but mixes pre-$\hpa$ spendable receipts into the counter; a consensus-level address is a consensus-level commitment and rotation is a hard fork absent a mechanism\\
O15 & emergency blocks and the drip & $D_\hgt = 0$ (required, S1); normal drip & otherwise 2-of-4 keys, facing no time gating, can accelerate the release of accumulated fees\\
O16 & forfeited tenant balances (\S23) & anywhere but $\pool$ (required, S2) & routing to the pool pays for report brigading\\
O17 & measured revenue baseline for $V_0$ & explorer receipts; list-price of measured resources & \S\ref{sec:pnr:crossover} scales linearly with it; closable by observation once $\mathrm{Rev}$ exists\\
O19 & tier partition basis & \PoN bases $1 : 3.5 : 9$ & confirm only; must not be described as collateral-proportional (Remark~\ref{rem:pnr:c24})\\
O20 & PA cut applied uniformly & uniform $\times 1/2$; other & changes operator issuance and every parity figure\\
\bottomrule
\end{tabular}
\caption{Parameters of the \PNR settlement mechanism, after the closures of
Table~\ref{tab:pnr:closed}; the values since settled are in \Cref{tab:pnr-recommended}, and O4, O8,
O14 and O19 remain open. The bold row is the one that gates the rest. Mirrored in
\texttt{plan/\allowbreak{}decisions-needed.\allowbreak{}md}.}
\label{tab:pnr:open}
\end{table}

Three requirements are exported to other mechanisms and are restated here so that they are not lost
at a section boundary: to \S16, \emph{every application-hosting charge, whatever the payment channel,
settles as exactly one L1 payment to $\mathsf{sink}$, or the consensus split is void and the revenue
counter under-reports}; to \S23,
\emph{forfeited balances never enter $\pool$} (Property~\ref{prop:pnr:exported}, S2); and to \S14,
\emph{the eligibility gate admitting a node to $\payees$ must be hosting-conditioned}
(Open Problem~\ref{op:pnr:d1}).

% ---------------------------------------------------------------------------
\subsection{Recommended values for the remaining open parameters}
\label{sec:pnr:recommended}

Each parameter still open above is given a recommended value here, with the reasoning that produced
it. Each was proposed by this paper and \textbf{confirmed by the author}
\srcintent{08-answers-round-2.md, round 11}, and several are simply a
recommendation to keep the default that \Cref{def:pnr:pool} already takes --- which is worth stating
explicitly, because a default that nobody has ratified is not the same object as a default that has
been examined and kept.

\begin{table}[H]
\centering\footnotesize
\caption{Values for the open \PNR{} parameters. Proposed by this paper and \textbf{confirmed by the
author} \srcintent{08-answers-round-2.md, round 11}. \planned, \srcinferred{} for the derivations.}
\label{tab:pnr-recommended}
\begin{tabular}{@{}llp{0.46\linewidth}@{}}
\toprule
& \textbf{Recommended} & Reason \\
\midrule
O5 tier with no confirmed node & \textbf{stays in $\pool$} & conservation-trivial; self-correcting
when that tier next confirms \\
O6 partition residue & \textbf{stays in $\pool$} & \num{0.021}\flux/year does not justify dev-fund
plumbing \\
O7 shielded payers & \textbf{moot} & pools retired; all settlement transparent
(\Cref{rem:o7-moot}) \\
O9 price rebasing & \textbf{rule-based} & removes a Foundation lever without adding an oracle
dependency \\
O10 drip ordering & \textbf{from $\pool_{\hgt-1}$} & keep the default: coinbase checkable before the
block's own receipts \\
O11 initial seed & \textbf{0} & a seed is a one-off Foundation transfer to operators \\
O12 fleet-growth guard & \textbf{review trigger at $2\times$} & documented, not automatic \\
O13 undershoot policy & \textbf{fixed schedule $+$ stated review} & no silent re-cutting \\
O15 emergency blocks & \textbf{$D_\hgt = 0$} & required by S1 \\
O17 revenue baseline & \textbf{trailing sink receipts} & measured, not assumed \\
O20 halving uniformity & \textbf{uniform $\times 1/2$} & every parity figure in
\S\ref{sec:pnr:crossover} assumes it \\
\bottomrule
\end{tabular}
\end{table}

\paragraph{The three that are not housekeeping.}

\textbf{O11, no seed.} Seeding $\pool_{\hpa-1}$ above zero would remove the ninety-day ramp and give
operators income from the first post-activation block. It would also be a one-off Foundation transfer
to node operators, and it would have to be described as one --- in a section whose entire argument is
that \PNR{} replaces issuance with demand-side revenue. Starting at zero makes the ramp visible and
makes the first ninety days an honest measurement of whether demand exists. If the ramp is judged too
harsh, the correct instrument is the activation schedule, not a subsidy that muddies what the pool
balance means.

\textbf{O9, rule-based rebasing.} The price table is today a Foundation lever over pool inflow with
no consensus involvement, which sits awkwardly beside a mechanism whose legitimacy rests on operators
being paid from receipts rather than from discretion. Of the three options, an oracle imports a trust
dependency that would need its own threat model and its own section; discretion preserves the lever.
A rule --- rebasing on a published function of trailing measured receipts, adjusted at fixed
intervals --- removes the lever without adding a dependency, and it is the only option under which a
tenant can predict next period's price from public data. What the rule should be is left open; that
it should be a rule is the recommendation.

\textbf{O17, trailing receipts rather than list-price value.} The placeholder baseline of
``\num{100}\flux{} per application per month'' should be replaced by the trailing receipts actually
observed at $\mathsf{sink}$, not by the list-price value of the measured resource totals. The two
differ by exactly the utilisation gap, and \S\ref{sec:limitations} records that realised utilisation
is not measurable from public data. A baseline computed from list prices would therefore embed an
unmeasured assumption in the one number the crossover analysis is most sensitive to, and would do so
invisibly. Trailing receipts are smaller, real, and reproducible by anyone watching the sink address.

\subsection{Comparison with prior art}
\label{sec:pnr:related}

\begin{table}[H]
\centering
\scriptsize
\begin{tabular}{lp{0.16\linewidth}p{0.16\linewidth}p{0.38\linewidth}}
\toprule
system & who is paid & from what & what \Flux would do differently\\
\midrule
Ethereum, post-EIP-1559 & block proposer (priority fee); base fee burned & per-transaction fee, same block & no burn; the operator share is pooled and dripped to \emph{all} nodes through the rotation, and a producer gains nothing by including a payment. The same-block capture EIP-1559 grants the proposer by design is precisely what Theorem~\ref{thm:pnr:timing} removes\\
Filecoin \citep{filecoin2017} & the specific storage provider & per-deal escrow released on proof of storage, with slashing & payment strictly conditional on proven service there; unconditional and network-wide here (H4, P7). \Flux would trade the proof-of-service dependency for fairness and simplicity, and inherit the free-rider problem of \S\ref{sec:pnr:hosting} in exchange\\
Akash \citep{akash} & the winning provider & tenant escrow drawn down per block at the bid price, with a take rate to a community pool & no per-lease escrow to the host and no auction; the network take would be \SIrange{80}{50}{\percent} to a Foundation address rather than a governance-controlled pool. Stated without softening\\
Render \citep{render} & node operators in proportion to work & burn-and-mint: payments burned, emissions minted to workers & burn-and-mint ties reward to metered work. The settled design rejects both the burn and the metering, which is exactly the roadmap model superseded by C1\\
Dash masternodes & one masternode per block, longest-since-paid queue & block subsidy plus a treasury share & the nearest structural precedent for Proposition~\ref{prop:pnr:rotation}; Dash never routed application fees through the rotation. \Flux would be that rotation plus a fee pool \srcdoc{plan/mech-pnr.md \S9}\\
Cosmos SDK \texttt{x/distribution} & all validators in proportion to stake, every block, lazily accumulated & transaction fees, with a community tax & exact and immediate, but requires per-validator accumulators and withdraw transactions. \Flux would pay through three existing coinbase outputs with no per-node state and no claim transaction (P4, P5), accepting a bounded \SI{3.83}{\percent} spread and a twenty-day lag instead. This is the road not taken \srcdoc{plan/mech-pnr.md \S9}\\
\bottomrule
\end{tabular}
\caption{\PNR against fee-handling and work-payment prior art.}
\label{tab:pnr:related}
\end{table}

The combination specific to \Flux is: coinbase-only payment with no per-node accounting; a
deterministic rotation that makes network-wide distribution provably fair over one rotation
(Propositions~\ref{prop:pnr:rotation} and~\ref{prop:pnr:within}); and a proportional drip with no
epoch boundary (Property~\ref{prop:pnr:invariants}, I3) that converts a schedule everyone can read
into a schedule nobody can profit from reading (Theorem~\ref{thm:pnr:timing}). The cost is stated
once more because it is the design's real weakness, and because \S\ref{sec:feasibility} and \S\ref{sec:migration} both depend on it:
\textbf{\PNR would pay for membership, and membership is only as real as the gate}.

%% Drain this section's float queue before the next begins.
\FloatBarrier
